\input _template

\setlanguage{SK}

\Title{Rovnaké dotyčnice}

\MathcompsLink{rovnake-dotycnice}

\Author{Patrik Bak}

\sec Úvod

V tomto letáku si ukážeme užitočnosť riešiteľskej techniky nazývanej \textit{rovnaké dotyčnice}. Opiera sa o~tie najzrejmejšie fakty, no úlohy, ktoré sa ňou dajú vyriešiť, môžu byť poriadne neintuitívne, a~preto sa budeme veľa venovať základom.

\sec Teória

\textit{Dohoda.} V celom tomto letáku, keď je $ABC$ trojuholník, píšeme $a = |BC|$, $b = |CA|$, $c = |AB|$ pre dĺžky jeho strán a~$s = \frac{a + b + c}{2}$ pre polovicu jeho obvodu.

\EnvId{BuRWesTIYiKA-BV2MAOvd-dve-dotycnice-z-bodu}
\Theorem{}{
    Nech je daná kružnica a~z bodu $A$ vedieme k~nej dve dotyčnice. Dotykové body označme $X$, $Y$. Potom $|AX| = |AY|$.

    \Image{equal-tangents-two-from-point.pdf}
}{
    Označme $S$ stred a~$r$ polomer danej kružnice. Z~bodu $A$ vedú dve rôzne dotyčnice, takže $A$ leží mimo kružnice a~je rôzny od oboch dotykových bodov. Dotyčnica je kolmá na polomer vedený do bodu dotyku, čiže $|\angle AXS| = |\angle AYS| = 90^\circ$ a~trojuholníky $AXS$, $AYS$ sú pravouhlé so spoločnou preponou $AS$ a~zhodnými odvesnami $|SX| = |SY| = r$. Pytagorova veta v~nich dáva
    $$
    |AX|^2 = |AS|^2 - r^2 = |AY|^2,
    $$
    a~keďže obe dĺžky sú nezáporné, je $|AX| = |AY|$.
}

\EnvId{mrvQMDl5Le3AmRPmyn1nx-spolocne-vonkajsie-dotycnice-dvoch-kruznic}
\Theorem{}{
    Nech $p$ a~$q$ sú spoločné vonkajšie dotyčnice dvoch kružníc $k_1$ a~$k_2$. Predpokladajme, že $p$ sa dotýka $k_1$ v~$A$ a~$k_2$ v~$B$ a~$q$ sa ich dotýka v~$C$ a~$D$. Potom $|AB| = |CD|$.

    \Image{equal-tangents-external-equal.pdf}
}{
    Predpokladajme najprv, že sa priamky $p$, $q$ pretínajú, a~ich priesečník označme $P$. Z~bodu $P$ vedú ku kružnici $k_1$ dotyčnice $p$ a~$q$ s~dotykovými bodmi $A$, $C$, takže podľa Tvrdenia 1 je $|PA| = |PC|$; rovnako pre $k_2$ s~dotykovými bodmi $B$, $D$ je $|PB| = |PD|$.

    Vonkajšia dotyčnica má obe kružnice v~tej istej polrovine, takže $k_1$ i~$k_2$ ležia v~tej istej z~oblastí, na ktoré dvojica priamok $p$, $q$ delí rovinu. Prienikom tejto oblasti s~priamkou $p$ je polpriamka so začiatkom $P$ a~ležia na nej aj body $A$, $B$, lebo tie ležia na $p$ i~na kružniciach. Body $A$, $B$ sú teda na tej istej polpriamke z~bodu $P$ a~podobne $C$, $D$ na tej istej polpriamke z~bodu $P$, odkiaľ
    $$
    |AB| = \bigl| |PA| - |PB| \bigr| = \bigl| |PC| - |PD| \bigr| = |CD|.
    $$

    Ak sa $p$, $q$ nepretínajú, sú rovnobežné. Kolmica na $p$ vedená stredom kružnice $k_1$ je kolmá aj na $q$ a~jej päty na $p$, $q$ sú práve dotykové body, čiže $A$, $C$ na nej ležia a~$AC \perp p$, pričom $|AC|$ je vzdialenosť priamok $p$, $q$. Rovnako $BD \perp p$. Posunutie, ktoré zobrazí $A$ na $C$, preto zobrazí priamku $p$ na $q$ a~bod $B$ na pätu kolmice z~$B$ na $q$, teda na $D$. Posunutie zachováva vzdialenosti, takže $|AB| = |CD|$.
}

\EnvId{jf0TF7EIF2sGDujPAaJGG-spolocna-vnutorna-dotycnica-dvoch-kruznic}
\Theorem{}{
    V predchádzajúcom nastavení predpokladajme naviac, že $k_1$ a~$k_2$ sú disjunktné, a~nech $r$ je spoločná vnútorná dotyčnica kružníc $k_1$ a~$k_2$, ktorá pretína $p$ v~$X$ a~$q$ v~$Y$. Potom $|XY| = |AB|$.

    \Image{equal-tangents-internal-equals-external.pdf}
}{
    Nech sa vnútorná dotyčnica $r$ dotýka kružnice $k_1$ v~bode $E$ a~kružnice $k_2$ v~bode $F$.

    Priamka $r$ je vnútorná dotyčnica, takže $k_1$ a~$k_2$ ležia v~opačných polrovinách s~hranicou $r$. Body $A \in k_1$ a~$B \in k_2$ sú teda priamkou $r$ oddelené a~úsečka $AB$ ju pretína; keďže $AB$ leží na $p$ a~priamky $p$, $r$ majú spoločný jediný bod $X$, leží $X$ na úsečke $AB$. Rovnako $Y$ leží na úsečke $CD$.

    Označme $H_p$, $H_q$ polroviny s~hranicami $p$, $q$ obsahujúce obe kružnice (tie existujú práve preto, že $p$, $q$ sú vonkajšie dotyčnice). Krajné body úsečky $AB$ ležia v~$H_q$, takže $X \in H_q$, a~rovnako $Y \in H_p$. Prienik $r \cap H_p$ je polpriamka so začiatkom $X$ a~obsahuje body $Y$, $E$, $F$; prienik $r \cap H_q$ je polpriamka so začiatkom $Y$ a~obsahuje body $X$, $E$, $F$. Dve polpriamky tej istej priamky, z~ktorých každá obsahuje začiatok tej druhej, majú za prienik úsečku spájajúcu ich začiatky, takže $E$ i~$F$ ležia na úsečke $XY$.

    Z~bodu $X$ vedú ku kružnici $k_1$ dotyčnice $p$ a~$r$ s~dotykovými bodmi $A$, $E$, takže podľa Tvrdenia 1 je $|XA| = |XE|$; rovnako $|XB| = |XF|$, $|YC| = |YE|$ a~$|YD| = |YF|$. Body $E$, $F$ ležia na úsečke $XY$, čiže $|XE| + |EY| = |XF| + |FY| = |XY|$, a~body $X$, $Y$ ležia na úsečkách $AB$, $CD$, čiže $|XA| + |XB| = |AB|$ a~$|YC| + |YD| = |CD|$. Spolu
    $$
    \gather
    2|XY| = (|XE| + |EY|) + (|XF| + |FY|) \\
    = (|XA| + |XB|) + (|YC| + |YD|) = |AB| + |CD|.
    \endgather
    $$
    Podľa Tvrdenia 2 je $|AB| = |CD|$, takže $2|XY| = 2|AB|$, teda $|XY| = |AB|$.
}

\EnvId{JhAKuB486cDd3AdwJYXkE-dotykove-body-vpisanej-kruznice}
\Theorem{}{
    Vpísaná kružnica trojuholníka $ABC$ sa dotýka strán $BC$, $CA$, $AB$ v~bodoch $D$, $E$, $F$. Potom $|AE| = |AF| = s - a = \frac{b + c - a}{2}$ a~podobne pre ostatné dva vrcholy.

    \Image{equal-tangents-incircle.pdf}
}{
    Vpísaná kružnica leží celá v~trojuholníku $ABC$, takže každý z~dotykových bodov leží aj v~trojuholníku; priesečníkom trojuholníka s~priamkou $BC$ je úsečka $BC$, čiže $D$ leží medzi $B$ a~$C$, a~rovnako $E$ leží medzi $C$ a~$A$ a~$F$ medzi $A$ a~$B$.

    Dve dotyčnice vedené ku kružnici z~toho istého bodu sú rovnako dlhé, takže môžeme položiť
    $$
    x = |AE| = |AF|, \qquad y = |BF| = |BD|, \qquad z = |CD| = |CE|.
    $$
    Z~uvedeného poradia bodov na stranách dostávame
    $$
    a = y + z, \qquad b = z + x, \qquad c = x + y.
    $$
    Sčítaním týchto troch rovností je $2(x + y + z) = a + b + c$, teda $x + y + z = s$, a~odčítaním prvej z~nich
    $$
    |AE| = |AF| = x = s - a = \frac{b + c - a}{2}.
    $$
    Rovnako $|BF| = |BD| = y = s - b$ a~$|CD| = |CE| = z = s - c$.
}

\EnvId{Ul1BcJ7hkSI2RgBb98Iik-polomer-vpisanej-kruznice-pravouhleho-trojuholnika}
\Theorem{}{
    Nech $ABC$ je pravouhlý trojuholník s~pravým uhlom pri vrchole $C$ a~nech $r$ je polomer jeho vpísanej kružnice. Potom
    $$
    r = s - c = \frac{a + b - c}{2}.
    $$

    \Image{equal-tangents-right-inradius.pdf}
}{
    Nech sa vpísaná kružnica so stredom $I$ dotýka strán $BC$, $CA$ v~bodoch $D$, $E$. Dotykový bod je pätou kolmice zo stredu, takže $ID \perp BC$ a~$IE \perp CA$; spolu s~pravým uhlom pri $C$ je teda $CDIE$ obdĺžnik. Jeho strany $ID$ a~$IE$ majú dĺžku $r$, ide teda o~štvorec a
    $$
    |CD| = |CE| = r.
    $$
    Podľa predchádzajúcej vety je zároveň $|CD| = |CE| = s - c$, čiže
    $$
    r = s - c = \frac{a + b - c}{2}.
    $$
}

\EnvId{4h53hlgz3wSLk4lVp4pyd-pripisana-kruznica-oproti-vrcholu}
\Theorem{}{
    Nech sa kružnica pripísaná trojuholníku $ABC$ oproti vrcholu $A$ dotýka priamok $BC$, $CA$, $AB$ v~bodoch $D$, $E$, $F$. Potom $|AE| = |AF| = s = \frac{a + b + c}{2}$.

    \Image{equal-tangents-excircle-from-a.pdf}
}{
    Najprv si ujasnime polohu dotykových bodov. Stred $I_A$ pripísanej kružnice leží na osi vnútorného uhla pri $A$, teda vnútri uhla $BAC$, a~od oboch priamok $AB$, $AC$ má vzdialenosť rovnú polomeru, takže celá kružnica leží v~uhle $BAC$ a~body $E$, $F$ ležia na polpriamkach $AC$, $AB$. Kružnica pripísaná oproti $A$ je zároveň zo svojej definície tá, ktorá leží v~polrovine s~hraničnou priamkou $BC$ neobsahujúcej $A$. Polpriamka $AC$ pretína priamku $BC$ jedine v~bode $C$, takže $E$ leží na predĺžení strany $AC$ za bodom $C$; rovnako $F$ leží na predĺžení strany $AB$ za bodom $B$. Bod $D$ leží na priamke $BC$ a~vnútri uhla $BAC$, čiže vnútri strany $BC$.

    Dve dotyčnice z~toho istého bodu sú rovnako dlhé, takže $|CE| = |CD|$, $|BF| = |BD|$ a~$|AE| = |AF|$. Z~poradia bodov je
    $$
    |AE| = b + |CD|, \qquad |AF| = c + |BD|, \qquad |BD| + |CD| = a.
    $$
    Sčítaním prvých dvoch rovností dostaneme $2|AE| = a + b + c$, teda
    $$
    |AE| = |AF| = s = \frac{a + b + c}{2}.
    $$
}

\EnvId{tIEpiTzploVkN1fSAAx_3-pripisana-kruznica-zvysne-vrcholy}
\Theorem{}{
    V predchádzajúcom nastavení platí $|BD| = |BF| = s - c = \frac{a + b - c}{2}$ a~$|CD| = |CE| = s - b = \frac{a + c - b}{2}$.

    \Image{equal-tangents-excircle-from-b-c.pdf}
}{
    V~dôkaze predchádzajúcej vety sme zistili polohu dotykových bodov: $F$ leží na polpriamke $AB$ za bodom $B$, bod $E$ na polpriamke $AC$ za bodom $C$ a~$D$ vnútri strany $BC$. Preto $|AF| = c + |BF|$ a~$|AE| = b + |CE|$, a~keďže $|AE| = |AF| = s$, dostávame
    $$
    |BF| = s - c = \frac{a + b - c}{2}, \qquad |CE| = s - b = \frac{a + c - b}{2}.
    $$
    Dotyčnice z~bodu $B$ dávajú $|BD| = |BF|$ a~dotyčnice z~bodu $C$ dávajú $|CD| = |CE|$, čo je dokazované tvrdenie.
}

\EnvId{G0QDSuOOrioc7nfucekbj-dotykove-body-symetricke-podla-stredu}
\Theorem{}{
    V predchádzajúcom nastavení sú body, v~ktorých sa vpísaná kružnica trojuholníka $ABC$ a~táto pripísaná kružnica dotýkajú priamky $BC$, symetrické podľa stredu strany $BC$.

    \Image{equal-tangents-excircle-symmetry.pdf}
}{
    Nech sa vpísaná kružnica trojuholníka $ABC$ dotýka strany $BC$ v~bode $P$ a~pripísaná kružnica oproti vrcholu $A$ sa dotýka priamky $BC$ v~bode $D$. Podľa vety o~dotykových bodoch vpísanej kružnice je $|BP| = s - b$ a~podľa predchádzajúcej vety $|BD| = s - c$, $|CD| = s - b$. Keďže $|BD| + |CD| = a = |BC|$, leží $D$ na úsečke $BC$; bod $P$ na nej leží tiež.

    Nech $M$ je stred strany $BC$. Stredová súmernosť podľa $M$ zobrazuje $B$ na $C$ a~úsečku $BC$ na seba, takže bod tejto úsečky vo vzdialenosti $s - b$ od $B$ zobrazí na bod tejto úsečky vo vzdialenosti $s - b$ od $C$. Prvým z~nich je $P$ a~druhým vďaka $|CD| = s - b$ bod $D$, čiže $P$ a~$D$ sú súmerné podľa $M$.
}

\EnvId{tER-vT4qkSmu9dw7bjhcr-dotycnicovy-stvoruholnik}
\Definition{Dotyčnicový štvoruholník}{
    Štvoruholník sa nazýva \textit{dotyčnicový}, ak má vpísanú kružnicu, teda kružnicu dotýkajúcu sa všetkých jeho štyroch strán.

    \Image{equal-tangents-pitot.pdf}
}

\EnvId{xFM0xDDSoTaKN7_k-PGs6-pitotova-veta}
\Theorem{Pitotova veta}{
    Konvexný štvoruholník $ABCD$ je dotyčnicový práve vtedy, keď $|AB| + |CD| = |BC| + |AD|$.
}{
    Nech má $ABCD$ vpísanú kružnicu, ktorá sa dotýka strán $AB$, $BC$, $CD$, $DA$ postupne v~bodoch $P$, $Q$, $R$, $S$. Z~rovnakých dotyčníc je $|AP| = |AS|$, $|BP| = |BQ|$, $|CQ| = |CR|$ a~$|DR| = |DS|$, takže
    $$
    \align
    |AB| + |CD| &= |AP| + |PB| + |CR| + |RD| \\
    &= |SA| + |BQ| + |QC| + |DS| \\
    &= \bigl(|BQ| + |QC|\bigr) + \bigl(|DS| + |SA|\bigr) \\
    &= |BC| + |AD|.
    \endalign
    $$

    Naopak nech konvexný štvoruholník $ABCD$ spĺňa $|AB| + |CD| = |BC| + |AD|$, čiže
    $$
    |AB| - |AD| = |BC| - |CD|.
    $$
    Zámena označení $B \leftrightarrow D$ zachováva konvexnosť aj túto podmienku, takže môžeme predpokladať $|AB| \geq |AD|$, a~teda aj $|BC| \geq |CD|$. Zvoľme bod $P$ na strane $AB$ s~$|AP| = |AD|$ a~bod $Q$ na strane $CB$ s~$|CQ| = |CD|$; potom
    $$
    |BP| = |AB| - |AD| = |BC| - |CD| = |BQ|.
    $$

    Označme $\alpha$, $\beta$ vnútorné uhly pri vrcholoch $A$, $B$. Ich vnútorné osi vychádzajú z~koncov úsečky $AB$ do tej istej polroviny a~zvierajú s~$AB$ uhly $\alpha/2$ a~$\beta/2$, obidva menšie než $90^\circ$, takže sa pretínajú v~jedinom bode $I$.

    Bod $D$ neleží na priamke $AB$, takže $P \neq D$ a~trojuholník $APD$ je rovnoramenný so základňou $PD$; os uhla pri $A$ je preto osou úsečky $PD$ a~$|IP| = |ID|$. Rovnako je os uhla pri $B$ osou úsečky $PQ$, takže $|IP| = |IQ|$ (v~prípade $P = Q$, ktorý nastane len pre $P = Q = B$, je to zrejmé). Spolu $|ID| = |IQ|$, čiže $I$ leží na osi úsečky $DQ$. Bod $D$ neleží ani na priamke $BC$, takže $D \neq Q$, a~keďže $|CD| = |CQ|$, je osou úsečky $DQ$ práve os uhla $DCB$.

    Bod $I$ má teda rovnakú vzdialenosť od priamok $AB$ a~$AD$, od priamok $AB$ a~$BC$ i~od priamok $CB$ a~$CD$, čiže od všetkých štyroch priamok $AB$, $BC$, $CD$, $DA$; označme ju $r$. Keby bolo $r = 0$, ležal by $I$ na priamke $AB$ aj na priamke $AD$, čiže $I = A$, no bod $A$ na osi uhla pri $B$ neleží. Preto $r > 0$ a~kružnica $\omega$ so stredom $I$ a~polomerom $r$ sa dotýka všetkých štyroch priamok.

    Zostáva overiť, že dotykové body ležia na stranách. Bod $I$ leží vnútri uhlov pri $A$ a~$B$, čiže vo vnútornej polrovine priamok $AB$, $AD$ aj $BC$. Na osi uhla pri $C$ leží $I$ na jej vnútornej polpriamke: opačná polpriamka totiž leží vo vrcholovom uhle, teda na opačnej strane priamky $BC$ než bod $A$. Bod $I$ tak leží vo vnútornej polrovine aj pre priamku $CD$, čiže vo vnútri štvoruholníka $ABCD$, a~preto aj vnútri uhla pri $D$; od priamok $DA$ a~$DC$ je rovnako vzdialený, takže leží na vnútornej osi uhla pri $D$.

    Nech je teraz $V$ ľubovoľný vrchol s~vnútorným uhlom $\varphi < 180^\circ$. Bod $I$ leží na vnútornej osi tohto uhla, takže s~ktorýmkoľvek jeho ramenom zviera pri $V$ uhol $\varphi/2 < 90^\circ$. Päta kolmice z~$I$ na toto rameno je preto od $V$ rôzna a~leží na ňom, nie na opačnej polpriamke: v~pravouhlom trojuholníku, ktorý vznikne, by opačná polpriamka dala pri $V$ tupý uhol. Dotykový bod kružnice $\omega$ s~priamkou $AB$ preto leží na polpriamke $AB$ aj na polpriamke $BA$, teda na strane $AB$, a~rovnako pre zvyšné tri strany. Kružnica $\omega$ je tak vpísanou kružnicou štvoruholníka $ABCD$.
}

\sec Úlohy

\EnvId{bm_8cGwtgRY12WAGCrh_a-pata-vysky-a-vpisane-kruznice}
\Problem{0}{DuoGeo 2025}{
    V trojuholníku $ABC$ nech $D$ je päta výšky z~vrcholu $C$, pričom $D$ leží na úsečke $AB$. Predpokladajme, že $|AC| + |BC|$ spolu s~priemermi vpísaných kružníc trojuholníkov $ADC$ a~$BDC$ dávajú v~súčte $26$ a~že $|CD| = 6$. Nájdite obsah trojuholníka $ABC$.
}{
    Jeden zo vzorcov, ktoré treba poznať, hovorí, ako vypočítať polomer vpísanej kružnice pravouhlého trojuholníka; v~našej úlohe máme ich dvojnásobky.
}{
    Ak je polomer vpísanej kružnice trojuholníka $ACD$ rovný $r_1$ a~polomer vpísanej kružnice trojuholníka $BCD$ je $r_2$, tak $2r_1=(|AD|+|CD|-|AC|)$ a~$2r_2=(|BD|+|CD|-|BC|)$. Ich sčítanie a~trochu úprav by nám mali dať posledný chýbajúci kúsok do výpočtu obsahu.
}{
    \Answer{Hľadaný obsah je $42$.}
    Označme $r_1$, $r_2$ polomery vpísaných kružníc trojuholníkov $ADC$, $BDC$; podmienka zo zadania hovorí, že $|AC| + |BC| + 2r_1 + 2r_2 = 26$. Keďže $CD$ je výška, sú oba tieto trojuholníky pravouhlé s~pravým uhlom pri $D$: ich odvesny sú $|AD|$, $|CD|$, resp. $|BD|$, $|CD|$, a~prepony $|AC|$, resp. $|BC|$. Vzorec pre polomer vpísanej kružnice pravouhlého trojuholníka preto dáva
    $$
    2r_1 = |AD| + |CD| - |AC|, \qquad 2r_2 = |BD| + |CD| - |BC|.
    $$

    Oba trojuholníky sú nedegenerované, takže $D$ je vnútorný bod úsečky $AB$ a~platí $|AD| + |BD| = |AB|$. Sčítaním oboch rovností tak dostaneme
    $$
    2r_1 + 2r_2 = |AB| + 2|CD| - |AC| - |BC|.
    $$

    Po dosadení do podmienky zo zadania sa členy $|AC|$ a~$|BC|$ vyrušia a~zostane
    $$
    26 = |AB| + 2|CD| = |AB| + 12,
    $$
    čiže $|AB| = 14$. Úsečka $CD$ je výška na stranu $AB$, a~preto
    $$
    S_{ABC} = \tfrac12 |AB| \cdot |CD| = \tfrac12 \cdot 14 \cdot 6 = 42.
    $$
}

\EnvId{EMMxpx8fQHdf5k7deQlkg-stred-strany-rovnobeznika}
\Problem{0}{MO 61-C-II-3}{
    Nech $E$ je stred strany $CD$ rovnobežníka $ABCD$, ktorý spĺňa $2\cdot|AB| = 3\cdot|BC|$. Dokážte, že ak má štvoruholník $ABCE$ vpísanú kružnicu, tak sa táto kružnica dotýka strany $BC$ v~jej strede.
}{
    Ak $|AB|=a$, tak $|BC|=2a/3$, $|CE|=a/2$. Vieme, že $ABCE$ je dotyčnicový. Čo vieme vyjadriť ďalej?
}{
    Mali by sme dostať $|AE|=|AB|+|CE|-|BC|=a+a/2-2a/3=5a/6$. Čo ešte vieme vyjadriť pomocou $a$?
}{
    Zjavne $|AD|=|BC|=2a/3$ a~$|DE|=|CE|=a/2$. Poznáme všetky strany trojuholníka $ADE$ pomocou $a$. Nevyzerajú podozrivo? Jeden pekný trik, ktorý je pri práci so zlomkami vo všeobecnosti užitočný, je zbaviť sa ich škálovaním -- označme $a=6x$. Aké sú teraz strany?
}{
    Mali by sme dostať, že $|AD|=4x$, $|DE|=3x$ a~$|AE|=5x$, čiže skrytý pravouhlý trojuholník so stranami 3-4-5. Takže náš rovnobežník je vlastne obdĺžnik. Teraz by už malo byť ľahké úlohu dokončiť.
}{
    Označme $a = |AB|$. Zo zadania je $|BC| = |AD| = \frac{2a}{3}$ a~z~rovnobežníka $|CD| = a$, takže pre stred $E$ strany $CD$ máme $|CE| = |DE| = \frac{a}{2}$.

    Bod $E$ leží vnútri strany $CD$, takže úsečky $AB$ a~$CE$ ležia na dvoch rôznych rovnobežných priamkach a~pri obehu $A \to B$, resp. $C \to E$ majú opačný smer; $ABCE$ je teda konvexný lichobežník so základňami $AB$ a~$CE$. Keďže má vpísanú kružnicu, dáva Pitotova veta $|AB| + |CE| = |BC| + |AE|$, čiže
    $$
    |AE| = a + \frac{a}{2} - \frac{2a}{3} = \frac{5a}{6}.
    $$

    Tým poznáme všetky tri strany trojuholníka $ADE$, sú to $\frac{2a}{3}$, $\frac{a}{2}$ a~$\frac{5a}{6}$. Zlomky vzťah medzi nimi zakrývajú, tak sa ich zbavme škálovaním: pri označení $a = 6x$ je
    $$
    |AD| = 4x, \qquad |DE| = 3x, \qquad |AE| = 5x,
    $$
    a~teda $|AD|^2 + |DE|^2 = 16x^2 + 9x^2 = 25x^2 = |AE|^2$. Podľa obrátenej Pytagorovej vety je uhol trojuholníka $ADE$ pri vrchole $D$ pravý.

    Bod $E$ leží vnútri úsečky $DC$, takže polpriamky $DE$ a~$DC$ splývajú a~$|\angle ADC| = |\angle ADE| = 90^\circ$. Rovnobežník s~jedným pravým uhlom je obdĺžnik, takže $ABCD$ je obdĺžnik; špeciálne $|\angle ABC| = 90^\circ$ a~$|\angle BCE| = |\angle BCD| = 90^\circ$.

    Nech $\omega$ je vpísaná kružnica štvoruholníka $ABCE$ so stredom $I$ a~polomerom $r$ a~nech sa strán $AB$, $BC$, $CE$ dotýka postupne v~bodoch $P$, $T$, $Q$. Platí $IP \perp AB$ a~$IT \perp BC$, a~keďže uhol pri $B$ je pravý, je $IP \parallel BC$ a~$IT \parallel AB$. Štvoruholník $BPIT$ je preto pravouholník a~$|BT| = |IP| = r$. Z~pravého uhla pri $C$ dostaneme rovnako $|CT| = |IQ| = r$.

    Bod $T$ leží na úsečke $BC$, takže $|BC| = |BT| + |CT| = 2r$ a~$|BT| = |CT|$. Kružnica $\omega$ sa teda dotýka strany $BC$ v~jej strede.
}

\EnvId{YHoydnnN3mLvYu3xOFGHp-obvod-trojuholnika-z-dotycnic}
\Problem{0}{Náboj 2008}{
    Nech $k$ je kružnica so stredom $S$ a~polomerom $1$ a~nech $P$ je bod s~$|PS| = 3$. Z~bodu $P$ vedieme dve dotyčnice ku $k$, ktoré sa jej dotýkajú v~bodoch $A$ a~$B$. Zvoľme ľubovoľný bod $T$ na kratšom oblúku $AB$ kružnice $k$ a~zostrojme dotyčnicu ku $k$ v~bode $T$. Táto dotyčnica pretína úsečky $AP$ a~$BP$ v~bodoch $X$ a~$Y$. Nájdite obvod trojuholníka $PXY$.
}{
    Naša kružnica má vzťah k~trojuholníku $PXY$, aký?
}{
    Je jeho kružnicou pripísanou oproti vrcholu $P$. To by nám malo dať dobrú predstavu, ako sa dá obvod $PXY$ interpretovať.
}{
    Pripomeňme si, že v~našich bežných konfiguráciách platí $|PA|=|PB|$, čo je polovica obvodu $PXY$. Takže stačí vypočítať $|PA|$.
}{
    \Answer{Obvod trojuholníka $PXY$ je $4\sqrt{2}$.}
    Celá úloha stojí na tom, čím je kružnica $k$ pre trojuholník $PXY$: je jeho kružnicou pripísanou oproti vrcholu $P$.

    Najprv poloha bodov. Ak by $T$ splynulo s~$A$ alebo s~$B$, dotyčnica v~$T$ by bola priamka $AP$, resp. $BP$, a~body $X$, $Y$ by neboli určené; bod $T$ je teda vnútorným bodom oblúka a~$T \neq A$, $T \neq B$. Priamka $XY$ sa dotýka $k$ v~$T$, takže s~$k$ nemá iný spoločný bod, a~preto na nej neležia body $A$ ani $B$.

    Kružnica $k$ sa dotýka priamok $PX$, $PY$, $XY$ postupne v~bodoch $A$, $B$, $T$. Dotykové body $A$, $B$ ležia na polpriamkach $PX$, $PY$ (dokonca za bodmi $X$, $Y$, lebo $X$ leží na úsečke $AP$ a~$Y$ na úsečke $BP$), takže $k$ je vpísaná do uhla $XPY$. Ďalej $k$ leží celá v~jednej z~polrovín určených priamkou $XY$, a~to v~tej, ktorá neobsahuje $P$: úsečka $PA$ pretína priamku $XY$ v~bode $X$, pričom ani $P$, ani $A$ na tejto priamke neležia. Kružnica dotýkajúca sa všetkých troch priamok $PX$, $PY$, $XY$, vpísaná do uhla $XPY$ a~ležiaca na opačnej strane priamky $XY$ než $P$, je práve kružnica pripísaná trojuholníku $PXY$ oproti vrcholu $P$.

    Podľa vety o~pripísanej kružnici oproti vrcholu sa vzdialenosť od $P$ k~jej dotykovým bodom na priamkach $PX$, $PY$ rovná polovici obvodu trojuholníka $PXY$, čiže
    $$
    |PA| = |PB| = \frac{|PX| + |XY| + |YP|}{2}.
    $$
    Stačí teda vypočítať $|PA|$. Dotyčnica je kolmá na polomer v~dotykovom bode, takže trojuholník $PAS$ má pri vrchole $A$ pravý uhol a~z~Pytagorovej vety
    $$
    |PA| = \sqrt{|PS|^2 - |SA|^2} = \sqrt{3^2 - 1^2} = 2\sqrt{2}.
    $$
    Obvod trojuholníka $PXY$ je preto $2|PA| = 4\sqrt{2}$, nezávisle od voľby bodu $T$.
}

\EnvId{njLBWjKeOLzDnXpHe7tic-vpisane-kruznice-na-uhlopriecke}
\Problem{0}{}{
    Nech $ABCD$ je dotyčnicový štvoruholník. Dokážte, že vpísané kružnice trojuholníkov $ABC$ a~$ADC$ sa dotýkajú uhlopriečky $AC$ v~tom istom bode.
}{
    Nech $X$ je dotykový bod vpísanej kružnice trojuholníka $ABC$ s~$AC$ a~$X'$ dotykový bod vpísanej kružnice trojuholníka $ADC$ s~$AC$. Stačí nám dokázať $|AX|=|AX'|$. Oba sú pekné výrazy.
}{
    Jediný trik je použiť vzorce typu $s-a$ v~trojuholníkoch $ABC$ a~$ADC$. Malo by nám zostať niečo, kde vieme použiť dotyčnicovosť.
}{
    Nech $X$ je bod, v~ktorom sa vpísaná kružnica trojuholníka $ABC$ dotýka strany $AC$, a~nech $X'$ je bod, v~ktorom sa tej istej úsečky dotýka vpísaná kružnica trojuholníka $ADC$. Oba ležia vo vnútri úsečky $AC$, teda na polpriamke $AC$, takže dokazované tvrdenie $X = X'$ je to isté ako $|AX| = |AX'|$. Obe tieto dĺžky sú pritom pekné, lebo pre ne máme vzorec typu $s - a$.

    Vzdialenosť vrcholu trojuholníka od dotykového bodu jeho vpísanej kružnice na priľahlej strane je polovica obvodu zmenšená o~protiľahlú stranu. V~trojuholníku $ABC$ je oproti vrcholu $A$ strana $BC$, čiže
    $$
    |AX| = \frac{|AB| + |AC| - |BC|}{2},
    $$
    a~v~trojuholníku $ADC$ je oproti $A$ strana $DC$, čiže
    $$
    |AX'| = \frac{|AD| + |AC| - |DC|}{2}.
    $$

    Člen $|AC|$ je v~oboch výrazoch ten istý, takže po odčítaní zostáva dokázať rovnosť $|AB| - |BC| = |AD| - |DC|$, čiže
    $$
    |AB| + |DC| = |BC| + |AD|.
    $$
    To však pre dotyčnicový štvoruholník $ABCD$ platí podľa Pitotovej vety, čím je dôkaz hotový.
}

\EnvId{6y2W6I3lf5ayFgDaZme17-tri-body-pripisanych-kruznic}
\Problem{0}{MO 63-B-I-3}{
    Na priamke $BC$ sú vyznačené tri body, v~ktorých sa pripísané kružnice trojuholníka $ABC$ dotýkajú tejto priamky; polohy samotných bodov $B$ a~$C$ nie sú dané. Nájdite bod, v~ktorom sa vpísaná kružnica trojuholníka $ABC$ dotýka priamky $BC$.
}{
    Nech sa kružnica pripísaná oproti vrcholu $C$ dotýka $BC$ v~$D$, pripísaná oproti $B$ v~$E$ a~pripísaná oproti $A$ v~$F$. Skús vyjadriť niektoré dĺžky pomocou našich vzorcov typu $s-a$.
}{
    Trik je uvedomiť si, že $|DB|=|CE|=s-a$. Čo to znamená? A~čo vieme o~$F$ a~o~dotykovom bode vpísanej kružnice?
}{
    \Answer{Spomedzi troch daných bodov leží práve jeden medzi zvyšnými dvomi (je to dotykový bod kružnice pripísanej oproti $A$). Hľadaný bod je jeho obrazom v~stredovej súmernosti podľa stredu úsečky určenej zvyšnými dvomi bodmi.}
    Nech sa kružnica pripísaná oproti vrcholu $C$ dotýka priamky $BC$ v~bode $D$, kružnica pripísaná oproti $B$ v~bode $E$ a~kružnica pripísaná oproti $A$ v~bode $F$; dotykový bod vpísanej kružnice s~priamkou $BC$, ktorý hľadáme, označme $T$. Body $B$, $C$ nepoznáme, takže musíme urobiť dve veci: rozoznať $D$, $E$, $F$ spomedzi troch daných bodov a~potom z~nich zostrojiť $T$.

    Dotyčnicový úsek z~vrcholu ku kružnici pripísanej oproti tomuto vrcholu má dĺžku $s$, takže $|CD| = s$ a~$|BE| = s$. Kružnica pripísaná oproti $C$ je vpísaná do uhla $ACB$, takže sa priamky $BC$ dotýka v~bode polpriamky $CB$; keďže $s - a = \frac{b + c - a}{2} > 0$, je $|CD| = s > a = |CB|$, čiže $D$ leží na tejto polpriamke za bodom $B$ a~
    $$
    |BD| = |CD| - |CB| = s - a.
    $$
    Rovnako $E$ leží na polpriamke $BC$ za bodom $C$ a~$|CE| = s - a$.

    Pre bod $F$ platí $|BF| = s - c$ a~$|CF| = s - b$; obe dĺžky sú kladné, takže $F$ je vnútorným bodom úsečky $BC$. Na priamke $BC$ teda ležia body v~poradí $D$, $B$, $F$, $C$, $E$. Bod $F$ je tak práve ten z~troch daných bodov, ktorý leží medzi zvyšnými dvomi, a~tým je rozoznaný; na rozlíšení $D$ od $E$ nezáleží, ako hneď uvidíme.

    Z~rovnosti $|BD| = |CE| = s - a$ vyplýva, že $D$ a~$E$ sú súmerné podľa stredu $M$ úsečky $BC$, inými slovami, $M$ je stredom úsečky $DE$. Podľa vety o~súmernosti dotykových bodov vpísanej kružnice a~kružnice pripísanej oproti $A$ na priamke $BC$ sú podľa toho istého bodu $M$ súmerné aj body $F$ a~$T$.

    Bod $T$ je preto obrazom bodu $F$ v~stredovej súmernosti podľa stredu úsečky $DE$, a~to už je konštrukcia využívajúca iba tri dané body. Ekvivalentne, $T$ je ten bod úsečky $DE$, pre ktorý
    $$
    |DT| = |EF| = (s - a) + (s - b).
    $$
}

\EnvId{wqOmXDet2xLsFXXqSqqXr-tri-dotycnice-vpisanej-kruznice}
\Problem{0}{}{
    Ku vpísanej kružnici trojuholníka sú zostrojené tri dotyčnice, každá z~nich odrezáva pri jednom z~vrcholov menší trojuholník. Predpokladajme, že tieto menšie trojuholníky majú obvody $1$, $2$ a~$3$. Dokážte, že pôvodný trojuholník je pravouhlý.
}{
    Nech je náš trojuholník $ABC$. Začnime jednou dotyčnicou, ktorá odrezáva trojuholník $AXY$. Ak sa naň pozrieme, čím je preň vpísaná kružnica trojuholníka $ABC$?
}{
    Kľúčové je uvedomiť si, že vpísaná kružnica je pripísanou kružnicou trojuholníka $AXY$ oproti vrcholu $A$. Čo nám to hovorí o~hodnote jeho obvodu?
}{
    To nám dáva, že polovica obvodu $AXY$ je vlastne vzdialenosť od $A$ k~dotykovým bodom vpísanej kružnice na $AB$, $AC$, teda $s-a$.
}{
    Posledným krokom je vyriešiť $s-a=1/2$, $s-b=1$, $s-c=3/2$ pre $a$, $b$, $c$.
}{
    Nech $ABC$ je náš trojuholník, $\omega$ jeho vpísaná kružnica a~nech sa $\omega$ dotýka strán $BC$, $CA$, $AB$ v~bodoch $D$, $E$, $F$.

    Ukážeme najprv, že trojuholník odrezaný pri vrchole $A$ má obvod $2(s-a)$ bez ohľadu na to, ktorú dotyčnicu sme zvolili. Nech teda dotyčnica ku kružnici $\omega$ pretína úsečku $AB$ v~bode $X$ a~úsečku $AC$ v~bode $Y$. Kružnica $\omega$ je vpísaná do uhla $XAY$, dotýka sa priamky $XY$ a~leží v~tej polrovine s~hraničnou priamkou $XY$, ktorá neobsahuje $A$. Je teda kružnicou pripísanou trojuholníku $AXY$ oproti vrcholu $A$. Úsečka $AF$ pritom pretína priamku $XY$, takže $X$ leží medzi $A$ a~$F$, a~analogicky $Y$ leží medzi $A$ a~$E$; body $F$ a~$E$ sú tak dotykovými bodmi tejto pripísanej kružnice s~predĺženiami strán $AX$ a~$AY$.

    Podľa vety o~pripísanej kružnici oproti vrcholu je $|AF| = |AE|$ rovné polovici obvodu trojuholníka $AXY$. Zároveň podľa vzorca pre dotykové body vpísanej kružnice je $|AF| = s - a$, a~teda obvod trojuholníka $AXY$ sa rovná $2(s-a)$.

    Obvod odrezaného trojuholníka teda závisí len od vrcholu, pri ktorom sme rezali. Označme vrcholy tak, aby pri $A$ bol odrezaný trojuholník s~obvodom $1$, pri $B$ s~obvodom $2$ a~pri $C$ s~obvodom $3$. Dostávame
    $$
    s - a = \tfrac12, \qquad s - b = 1, \qquad s - c = \tfrac32.
    $$
    Súčet ľavých strán je $3s - (a+b+c) = s$, takže $s = \tfrac12 + 1 + \tfrac32 = 3$ a~odtiaľ
    $$
    a = \tfrac52, \qquad b = 2, \qquad c = \tfrac32.
    $$
    Takéto strany naozaj tvoria trojuholník, lebo $b + c = \tfrac72 > a$. Napokon
    $$
    b^2 + c^2 = 4 + \tfrac94 = \tfrac{25}{4} = a^2,
    $$
    takže podľa obrátenej Pytagorovej vety má trojuholník $ABC$ pravý uhol pri vrchole $A$, teda pri tom vrchole, kde je odrezaný trojuholník s~obvodom $1$.
}

\EnvId{ywat5ogtaZvQd-symWXYX-dotykove-body-na-uhloprieckach-rovnobeznika}
\Problem{0}{MO 54-A-S-2}{
    Nech $ABCD$ je rovnobežník s~$|AB| > |BC|$. Nech $K$ a~$M$ sú body, v~ktorých sa vpísané kružnice trojuholníkov $ACD$ a~$ABC$ dotýkajú uhlopriečky $AC$, a~nech $L$ a~$N$ sú analogicky definované pre trojuholníky $BCD$ a~$ABD$ dotýkajúce sa uhlopriečky $BD$. Dokážte, že $KLMN$ je obdĺžnik.
}{
    Rovnobežník je symetrický podľa priesečníka uhlopriečok. Čo to znamená pre $KLMN$?
}{
    Pri symetrii rovnobežníka sa zobrazia aj vpísané kružnice, a~tak aj ich dotykové body. Teda ak $P$ je priesečník uhlopriečok, tak $|PK|=|PM|$ a~$|PN|=|PL|$. Čo to znamená a~čo zostáva ukázať?
}{
    Máme, že $KLMN$ je rovnobežník. Aby sme ukázali, že je to naozaj obdĺžnik, stačí ukázať, že jeho uhlopriečky sú rovnaké.
}{
    Pozrime sa na $KM$. Rovná sa $|AC|-|AK|-|CM|$. Vďaka symetrii je to $|AC|-2|AK|$. Ale $|AK|$ vieme vyjadriť pomocou našich vzorcov typu $s-a$.
}{
    Nech $P$ je priesečník uhlopriečok a~nech $\sigma$ je stredová súmernosť so stredom $P$. Uhlopriečky rovnobežníka sa rozpoľujú, takže $\sigma(A) = C$, $\sigma(B) = D$, $\sigma(C) = A$, $\sigma(D) = B$, a~teda $\sigma$ zobrazuje trojuholník $ACD$ na trojuholník $CAB$ a~trojuholník $BCD$ na trojuholník $DAB$. Je to zhodnosť, preto zobrazí vpísanú kružnicu na vpísanú kružnicu; priamku $AC$ aj priamku $BD$ pritom zobrazí samu na seba, takže dotykový bod prejde na dotykový bod. Dostávame $\sigma(K) = M$ a~$\sigma(L) = N$, čiže $P$ je stredom oboch úsečiek $KM$ a~$LN$.

    Úsečky $KM$ a~$LN$ sa teda navzájom rozpoľujú a~ležia na rôznych priamkach $AC$, $BD$. Ak ukážeme $|KM| = |LN| > 0$, budú všetky štyri body rôzne a~$KLMN$ bude rovnobežník s~rovnako dlhými uhlopriečkami, teda obdĺžnik: body $K$, $L$, $M$, $N$ vtedy ležia na kružnici so stredom $P$ a~polomerom $|KM|/2$, úsečky $KM$, $LN$ sú jej rôzne priemery a~Tálesova veta dáva pri každom z~vrcholov $K$, $L$, $M$, $N$ pravý uhol.

    Označme $x = |AB| - |BC| > 0$ a~využime, že v~rovnobežníku platí $|AD| = |BC|$ a~$|CD| = |AB|$. Bod $K$ je dotykovým bodom vpísanej kružnice trojuholníka $ACD$ so stranou $AC$ a~bod $M$ dotykovým bodom vpísanej kružnice trojuholníka $ABC$ s~tou istou stranou, takže podľa vzorca $s-a$ pre dotykové body vpísanej kružnice
    $$
    \gather
    |AK| = \frac{|AC| + |AD| - |CD|}{2} = \frac{|AC| - x}{2}, \\
    |CM| = \frac{|CA| + |CB| - |AB|}{2} = \frac{|AC| - x}{2}.
    \endgather
    $$
    Obe dĺžky sú menšie než $|AC|/2$, preto body idú na úsečke $AC$ v~poradí $A$, $K$, $M$, $C$ a
    $$
    |KM| = |AC| - |AK| - |CM| = x.
    $$

    Rovnako je $L$ dotykovým bodom vpísanej kružnice trojuholníka $BCD$ so stranou $BD$ a~$N$ dotykovým bodom vpísanej kružnice trojuholníka $ABD$ s~ňou, takže
    $$
    \gather
    |BL| = \frac{|BD| + |BC| - |CD|}{2} = \frac{|BD| - x}{2}, \\
    |DN| = \frac{|DB| + |DA| - |AB|}{2} = \frac{|BD| - x}{2},
    \endgather
    $$
    a~z~rovnakého dôvodu $|LN| = |BD| - |BL| - |DN| = x$.

    Teda $|KM| = |LN| = x > 0$ a~$KLMN$ je obdĺžnik.
}

\EnvId{1rKRGiDDH6a1N3ui7-NIp-dotykove-body-v-rozdelenom-trojuholniku}
\Problem{0}{MO 64-A-I-5}{
    V trojuholníku $ABC$ nech $D$ je bod, v~ktorom sa vpísaná kružnica dotýka strany $BC$. Vpísaná kružnica trojuholníka $ABD$ sa dotýka strán $AB$ a~$BD$ v~bodoch $K$ a~$L$ a~vpísaná kružnica trojuholníka $ADC$ sa dotýka strán $DC$ a~$AC$ v~bodoch $M$ a~$N$. Dokážte, že body $K$, $L$, $M$, $N$ ležia na jednej kružnici.
}{
    Ak by naše body ležali na kružnici, čo by bolo jej stredom?
}{
    Pozri sa na osi úsečiek $KL$ a~$MN$. Čo sú a~kde sa pretínajú?
}{
    Vďaka $|BK|=|BL|$ a~$|CM|=|CN|$ musia byť osami uhlov pri $B$ a~$C$ v~trojuholníku $ABC$, takže by sa mali pretínať v~$I$. Preto zostáva dokázať, že $|IM|=|IL|$.
}{
    Keďže $D$ je dotykový bod vpísanej kružnice, platí $ID \perp BC$. Čo potrebujeme dokázať, aby sme dokázali $|IM|=|IL|$?
}{
    Posledným krokom je dokázať $|DM|=|DL|$. Ale to bude ľahké s~naším vzorcom typu $s-a$.
}{
    Stačí ukázať $(|BD|+|AD|-|AB|)/2=(|CD|+|AD|-|AC|)/2$. Ale $|BD|$, $|CD|$ vieme ľahko vypočítať pomocou $a$, $b$, $c$.
}{
    Označme $I$ stred vpísanej kružnice trojuholníka $ABC$ a~$r$ jej polomer. Stred hľadanej kružnice musí ležať na osi úsečky $KL$ aj na osi úsečky $MN$, a~práve tieto dve osi vieme identifikovať: sú to osi uhlov pri $B$ a~$C$, takže stredom bude $I$.

    Najprv poloha bodu $D$. Podľa vzorcov pre dotykové body vpísanej kružnice je $|BD| = s - b$ a~$|CD| = s - c$, čo sú obe kladné čísla, takže $D$ je vnútorný bod strany $BC$. Trojuholníky $ABD$ a~$ADC$ sú teda nedegenerované, polpriamka $BD$ je polpriamkou $BC$ a~polpriamka $CD$ je polpriamkou $CB$. Špeciálne $L$ leží na polpriamke $BC$ a~$M$ na polpriamke $CB$.

    Vpísaná kružnica trojuholníka $ABD$ sa dotýka strán $BA$ a~$BD$ v~bodoch $K$ a~$L$, takže $|BK| = |BL|$. Osová súmernosť podľa osi uhla $ABC$ zobrazuje polpriamku $BA$ na polpriamku $BC$, a~keďže $K$, $L$ ležia na týchto polpriamkach v~rovnakej vzdialenosti od $B$, zobrazuje $K$ na $L$. Body $K$, $L$ sú rôzne: platí $|BK| = (|AB| + |BD| - |AD|)/2 > 0$ podľa trojuholníkovej nerovnosti v~$ABD$, takže ani jeden z~nich nesplýva s~$B$, a~pritom ležia na dvoch rôznych polpriamkach z~$B$. Os uhla $ABC$ je teda osou úsečky $KL$. Rovnako $|CM| = |CN|$ a~súmernosť podľa osi uhla $ACB$ zobrazuje $N$ na $M$, čiže os uhla $ACB$ je osou úsečky $MN$.

    Obe osi uhlov prechádzajú bodom $I$, a~preto
    $$
    |IK| = |IL|, \qquad |IM| = |IN|.
    $$
    Zostáva dokázať $|IL| = |IM|$.

    Tu využijeme, že $D$ je dotykový bod vpísanej kružnice trojuholníka $ABC$ so stranou $BC$: platí $ID \perp BC$ a~$|ID| = r$. Body $L$, $M$ ležia na priamke $BC$, takže z~Pytagorovej vety
    $$
    |IL|^2 = r^2 + |DL|^2, \qquad |IM|^2 = r^2 + |DM|^2,
    $$
    a~dokazovaná rovnosť je ekvivalentná s~$|DL| = |DM|$.

    Obe tieto dĺžky sú dotyčnicové úseky z~vrcholu $D$, stačí teda použiť vzorec typu $s - a$ v~trojuholníku $ABD$ a~v~trojuholníku $ADC$:
    $$
    \gather
    |DL| = \frac{|BD| + |AD| - |AB|}{2}, \\
    |DM| = \frac{|CD| + |AD| - |AC|}{2}.
    \endgather
    $$
    Člen $|AD|$ sa v~rozdiele vyruší, takže $|DL| = |DM|$ je to isté ako
    $$
    |BD| - |CD| = |AB| - |AC| = c - b.
    $$
    A~to už vieme: $|BD| - |CD| = (s - b) - (s - c) = c - b$.

    Dostávame $|IK| = |IL| = |IM| = |IN|$, pričom táto spoločná hodnota je aspoň $|ID| = r > 0$. Body $K$, $L$, $M$, $N$ teda ležia na kružnici so stredom $I$ a~polomerom $|IL|$.
}

\EnvId{XY8K2xIEBBfNKdV7OHBv4-pripisana-kruznica-stvoruholnika}
\Problem{0}{}{
    Hovoríme, že konvexný štvoruholník $ABCD$ má pripísanú kružnicu oproti vrcholu $A$, ak je vpísaná do uhla $BAD$, dotýka sa priamok $BC$ a~$CD$ a~leží mimo štvoruholníka. Dokážte, že pre taký štvoruholník platí $|AB| + |CB| = |AD| + |CD|$.
}{
    Dotykové body kružnice s~$AB$, $BC$, $CD$, $DA$ označme $A'$, $B'$, $C'$, $D'$. Skús vyhodnotiť $|AB|+|BC|$ pomocou rovnakých dotyčníc.
}{
    Výpočet začína $|AB|+|BC|=|AA'|-|BA'|+|BC|$.
}{
    Nechajme $|AA'|$ nedotknuté a~urobme niečo s~$|BA'|$, čo nám pomôže využiť člen $+|BC|$.
}{
    Ďalší krok je $|BA'|=|BB'|=|BC|+|CB'|$. To dáva $|AB|+|BC|=|AA'|-|CB'|$. Teraz by $|AD|+|CD|$ malo dať niečo podobné.
}{
    Označme $\omega$ kružnicu zo zadania a~jej dotykové body s~priamkami $AB$, $BC$, $CD$, $DA$ postupne $A'$, $B'$, $C'$, $D'$. Keďže $\omega$ je vpísaná do uhla $BAD$, leží bod $A'$ na polpriamke $AB$ a~bod $D'$ na polpriamke $AD$.

    Žiadny z~vrcholov $A$, $B$, $C$, $D$ neleží na $\omega$: každý z~nich je priesečníkom dvoch z~priamok $AB$, $BC$, $CD$, $DA$, a~keby ležal na $\omega$, boli by tieto dve rôzne priamky dotyčnicami $\omega$ v~tom istom bode. Z~toho istého dôvodu sú body $A'$, $B'$, $C'$, $D'$ navzájom rôzne. Každý vrchol teda leží mimo $\omega$ a~dĺžky dotyčníc z~neho sú rovnaké:
    $$
    |AA'| = |AD'|, \quad |BA'| = |BB'|, \quad |CB'| = |CC'|, \quad |DC'| = |DD'|.
    $$

    Celý výpočet stojí na polohe týchto dotykových bodov, tak si ju najprv ujasníme.

    Bod $B$ leží medzi $A$ a~$A'$. Kružnica $\omega$ je totiž vpísaná do uhla $BAD$, takže leží v~polrovine určenej priamkou $AB$ a~obsahujúcej $D$, a~priamky $AB$ sa dotýka v~$A'$; všetky jej body okrem $A'$ preto ležia v~otvorenej polrovine, teda mimo priamky $AB$. Keby bol $A'$ vnútorným bodom strany $AB$, ležali by body $\omega$ dostatočne blízko k~$A'$ vo vnútri konvexného štvoruholníka $ABCD$, čo odporuje tomu, že $\omega$ leží mimo neho. Bod $A'$ leží na polpriamke $AB$ a~nesplýva ani s~$A$, ani s~$B$, takže na nej leží až za bodom $B$. Rovnako $D$ leží medzi $A$ a~$D'$.

    Odtiaľ dostaneme polohu celej $\omega$. Úsečka $AA'$ pretína priamku $BC$ v~bode $B$ a~ani jeden z~bodov $A$, $A'$ na tejto priamke neleží, takže $A$ a~$A'$ ležia v~opačných polrovinách určených priamkou $BC$. Kružnica $\omega$ sa priamky $BC$ dotýka, čiže celá leží v~jednej z~týchto polrovín, a~keďže $A'$ je jej bodom, je to tá, ktorá neobsahuje $A$. Z~toho, že $D$ leží medzi $A$ a~$D'$, rovnako dostaneme, že $\omega$ leží v~polrovine určenej priamkou $CD$ a~neobsahujúcej $A$.

    Teraz vieme umiestniť aj $B'$ a~$C'$. Z~konvexnosti ležia body $A$ a~$B$ v~tej istej polrovine určenej priamkou $CD$, kým $B'$ leží v~tej opačnej, a~to mimo priamky $CD$ (inak by $B'$ bol spoločným bodom $\omega$ a~priamky $CD$, čiže $B' = C'$). Úsečka $BB'$ preto pretína priamku $CD$; leží však na priamke $BC$, ktorá má s~priamkou $CD$ jediný spoločný bod $C$. Bod $C$ tak leží medzi $B$ a~$B'$ a~rovnaký argument so zámenou $B \leftrightarrow D$ dáva, že $C$ leží medzi $D$ a~$C'$.

    Zvyšok je výpočet. Keďže $B$ leží medzi $A$ a~$A'$ a~bod $C$ medzi $B$ a~$B'$, platí
    $$
    \eqalign{
    |AB| + |CB| &= |AA'| - |BA'| + |BC| = \cr
    &= |AA'| - |BB'| + |BC| = \cr
    &= |AA'| - \bigl(|BC| + |CB'|\bigr) + |BC| = \cr
    &= |AA'| - |CB'|. \cr
    }
    $$
    Úplne rovnako, s~$D$ medzi $A$ a~$D'$ a~s~bodom $C$ medzi $D$ a~$C'$, dostaneme
    $$
    |AD| + |CD| = |AD'| - |CC'|.
    $$
    Pravé strany sa rovnajú, lebo $|AA'| = |AD'|$ a~$|CB'| = |CC'|$.
}

\EnvId{4CAVBeHRTqtsMdY8jrKuX-vpisane-kruznice-dotykajuce-sa-zvonka}
\Problem{0}{MO 49-A-I-2}{
    Nech $K$, $L$, $M$ sú vnútorné body strán $BC$, $CA$, $AB$ trojuholníka $ABC$, zvolené tak, že vpísané kružnice trojuholníkov $ABK$ a~$CAK$ sa dotýkajú zvonka, vpísané kružnice $BCL$ a~$ABL$ sa dotýkajú zvonka a~vpísané kružnice $CAM$ a~$BCM$ sa dotýkajú zvonka. Dokážte, že
    $$
    |BK| \cdot |CL| \cdot |AM| = |CK| \cdot |AL| \cdot |BM|.
    $$
}{
    Najprv zabudni na $L$ a~$M$ a~nakresli len obrázok s~$K$. Dĺžku $|BK|$ vieme vyjadriť pomocou $a$, $b$, $c$ a~to bude všetko.
}{
    Nech~$D$ je spoločný dotykový bod. Dĺžku $|AD|$ vieme vyjadriť dvomi spôsobmi pomocou vzorca typu $s-a$ v~$ABK$ aj v~$AKC$. Porovnaj tieto výrazy.
}{
    Po ich porovnaní by sme mali dostať $|AD|=(c+|AK|-|BK|)/2$ z~$ABK$ a~$|AD|=(b+|AK|-|CK|)/2$ z~toho druhého. Polož ich do rovnosti a~použi aj $|BK|+|CK|$. Čo to dá pre $|BK|$?
}{
    Zabudnime nateraz na body $L$ a~$M$ a~vypočítajme $|BK|$ pomocou $a$, $b$, $c$; tým bude úloha hotová.

    Bod $K$ leží vnútri strany $BC$, takže body $B$ a~$C$ ležia v~opačných polrovinách s~hraničnou priamkou $AK$. Vpísaná kružnica trojuholníka $ABK$ leží v~prvej z~nich, vpísaná kružnica trojuholníka $CAK$ v~druhej a~obe sa dotýkajú priamky $AK$. Podľa predpokladu majú tieto dve kružnice spoločný bod; ten musí ležať v~priesečníku oboch polrovín, čiže na priamke $AK$. Každá z~kružníc má však s~$AK$ jediný spoločný bod, a~to svoj dotykový bod, takže tieto dva dotykové body splývajú. Označme ich $D$.

    Vzdialenosť od vrcholu k~dotykovému bodu vpísanej kružnice na susednej strane je vzorec typu $s-a$, takže z~trojuholníka $ABK$ dostávame
    $$
    |AD| = \frac{c + |AK| - |BK|}{2}
    $$
    a~z trojuholníka $CAK$
    $$
    |AD| = \frac{b + |AK| - |CK|}{2}.
    $$
    Porovnaním sa $|AK|$ vyruší a~zostáva $c - |BK| = b - |CK|$, čiže $|CK| - |BK| = b - c$. Keďže $K$ leží vnútri $BC$, platí navyše $|BK| + |CK| = a$, a~teda
    $$
    \gather
    |BK| = \frac{a + c - b}{2} = s - b, \\
    |CK| = \frac{a + b - c}{2} = s - c.
    \endgather
    $$

    Zadanie je invariantné voči cyklickej zámene $A \to B \to C \to A$, pri ktorej $K \to L \to M$: strana $BC$ prejde na $CA$, trojuholník $ABK$ na $BCL$ a~trojuholník $CAK$ na $ABL$, čo je presne podmienka kladená na $L$; ešte jedno použitie dá podmienku na $M$. Pri tejto zámene sa $a \to b \to c \to a$, takže $s$ zostáva rovnaké. Predchádzajúci výpočet preto rovno dáva
    $$
    \gather
    |CL| = s - c, \qquad |AL| = s - a, \\
    |AM| = s - a, \qquad |BM| = s - b.
    \endgather
    $$
    Obe strany dokazovanej rovnosti sú tak tým istým súčinom:
    $$
    \eqalign{
    |BK| \cdot |CL| \cdot |AM| &= (s-a)(s-b)(s-c) = \cr
    &= |CK| \cdot |AL| \cdot |BM|. \cr
    }
    $$
}

\EnvId{mm753DtWsC9P9vCKmpLmt-devat-mensich-stvoruholnikov}
\Problem{0}{}{
    Na každej strane konvexného štvoruholníka sú zvolené dva body a~dvojice protiľahlých bodov sú spojené úsečkami tak, že sa spájajúce úsečky nekrížia. Pôvodný štvoruholník je tým rozdelený na deväť menších štvoruholníkov. Ukážte, že ak má každý zo štyroch rohových štvoruholníkov a~stredný štvoruholník vpísanú kružnicu, tak ju má aj pôvodný štvoruholník.

    \Image{equal-tangents-checkerboard.pdf}
}{
    Potrebujeme dokázať $|AB|+|CD|=|BC|+|AD|$. Rohové kružnice označme $w_a, w_b, w_c, w_d$ (pomenované podľa najbližšieho vrcholu $ABCD$) a~strednú kružnicu $\omega$. Skús vyjadriť každú stranu $ABCD$ pomocou jednoduchších častí.
}{
    Strana $AB$ obsahuje dotyčnicový úsek z~$A$ ku $w_a$, dotyčnicový úsek z~$B$ ku $w_b$ a~časť medzi nimi. Tá celá stredná časť je len spoločná vonkajšia dotyčnica $w_a$ a~$w_b$ pozdĺž priamky $AB$. Rozlož zvyšné tri strany rovnako a~pomocou rovnakých dotyčníc zredukuj, čo potrebujeme dokázať.
}{
    Trik je v~tom, že dĺžka dotyčnice z~$A$ ku $w_a$ je súčasťou $AB$ aj $AD$, a~podobne pre ostatné rohy. Takže stačí dokázať niečo o~vonkajších dotyčniciach.
}{
    Nech $d(x,y)$ označuje dĺžku vonkajšej dotyčnice dvoch rohových kružníc pozdĺž tej strany $ABCD$, ktorej sa obe dotýkajú. Zostáva nám dokázať $d(w_a,w_b)+d(w_c,w_d)=d(w_b,w_c)+d(w_a,w_d)$. Aby sme to dosiahli, musíme tieto dĺžky nájsť inde než na stranách $ABCD$. Tu prichádza na scénu $\omega$.
}{
    Každá dvojica susediacich rohových kružníc sa dotýka aj jednej zo štyroch vnútorných úsečiek (tej, ktorá ohraničuje okrajovú bunku zdieľanú týmito dvomi rohmi voči strednej bunke). Preneste každú zo štyroch vonkajších dotyčníc na zodpovedajúcu vnútornú úsečku pomocou rovnakých dotyčníc.
}{
    Na každej vnútornej úsečke teraz prenesená dotyčnica prechádza cez dva vrcholy strednej bunky, kde sa dotýka aj $\omega$. Rovnaké dotyčnice z~týchto vrcholov rozdelia každý prenesený kúsok na dve polovice, ktoré sa zhodujú s~polovicami na susediacich vnútorných úsečkách, presne tak, ako sa rohové dotyčnice zhodovali v~$A$, $B$, $C$, $D$ na vonkajšej hranici. Tie štyri kúsky sa potom po dvojiciach vyrušia a~sme hotoví.
}{
    Zvolené body označme podľa poradia na obvode: nech $P_1$, $P_2$ ležia na $AB$ v~poradí $A$, $P_1$, $P_2$, $B$, body $Q_1$, $Q_2$ na $BC$ v~poradí $B$, $Q_1$, $Q_2$, $C$, body $R_1$, $R_2$ na $CD$ v~poradí $C$, $R_1$, $R_2$, $D$ a~body $S_1$, $S_2$ na $DA$ v~poradí $D$, $S_1$, $S_2$, $A$. Spájajú sa protiľahlé strany, čiže $AB$ s~$CD$ a~$BC$ s~$DA$, a~podmienka, že sa spájajúce úsečky nekrížia, dáva práve dvojice $P_1R_2$, $P_2R_1$ a~$Q_1S_2$, $Q_2S_1$. Vrcholy strednej bunky sú
    $$
    \gather
    E = P_1R_2 \cap Q_1S_2, \qquad F = P_2R_1 \cap Q_1S_2, \\
    G = P_2R_1 \cap Q_2S_1, \qquad H = P_1R_2 \cap Q_2S_1,
    \endgather
    $$
    takže deväť buniek je toto: rohové $AP_1ES_2$, $P_2BQ_1F$, $Q_2CR_1G$, $R_2DS_1H$ (postupne pri vrcholoch $A$, $B$, $C$, $D$), okrajové $P_1P_2FE$, $Q_1Q_2GF$, $R_1R_2HG$, $S_1S_2EH$ a~stredná $EFGH$. Vpísané kružnice rohových buniek označme $w_a$, $w_b$, $w_c$, $w_d$ podľa najbližšieho vrcholu a~vpísanú kružnicu strednej bunky $\omega$.

    Štvoruholník $ABCD$ je konvexný, takže podľa Pitotovej vety stačí dokázať
    $$
    |AB| + |CD| = |BC| + |AD|. \eqno(*)
    $$

    Kružnice $w_a$ a~$w_b$ sa obe dotýkajú priamky $AB$, prvá vo vnútornom bode $T_A$ úsečky $AP_1$ a~druhá vo vnútornom bode $T_B$ úsečky $P_2B$ (to sú strany ich buniek ležiace na $AB$). Označme $d_{AB} = |T_AT_B|$ a~analogicky $d_{BC}$, $d_{CD}$, $d_{AD}$ pre zvyšné tri strany a~príslušné dvojice susedných rohových kružníc. Kružnica $w_a$ sa dotýka aj priamky $AD$, takže podľa vety o~dvoch dotyčniciach z~bodu sa $|AT_A|$ rovná vzdialenosti bodu $A$ od dotykového bodu na $AD$; túto spoločnú hodnotu označme $t_A$ a~podobne zaveďme $t_B$, $t_C$, $t_D$. Z~poradia $A$, $T_A$, $T_B$, $B$ na strane $AB$ a~z~rovnakých rozkladov zvyšných strán dostávame
    $$
    \eqalign{
    |AB| &= t_A + d_{AB} + t_B, \cr
    |BC| &= t_B + d_{BC} + t_C, \cr
    |CD| &= t_C + d_{CD} + t_D, \cr
    |AD| &= t_A + d_{AD} + t_D. \cr
    }
    $$
    Každá zo štyroch dotyčnicových dĺžok $t_A$, $t_B$, $t_C$, $t_D$ vystupuje raz na ľavej a~raz na pravej strane $(*)$, takže $(*)$ je ekvivalentné s~rovnosťou
    $$
    d_{AB} + d_{CD} = d_{BC} + d_{AD}. \eqno(**)
    $$

    Kľúčový krok je prečítať dĺžky $d$ inde než na stranách $ABCD$. Kružnice $w_a$, $w_b$ sa totiž okrem $AB$ dotýkajú aj priamky $Q_1S_2$: prvá ako vpísaná kružnica bunky $AP_1ES_2$ vo vnútornom bode $U_A$ úsečky $S_2E$, druhá ako vpísaná kružnica bunky $P_2BQ_1F$ vo vnútornom bode $U_B$ úsečky $FQ_1$. Obe tieto bunky ležia v~tej polrovine s~hraničnou priamkou $Q_1S_2$, ktorá obsahuje stranu $AB$, čiže $AB$ aj $Q_1S_2$ sú spoločné vonkajšie dotyčnice kružníc $w_a$, $w_b$. Podľa vety o~spoločných vonkajších dotyčniciach dvoch kružníc je preto
    $$
    d_{AB} = |T_AT_B| = |U_AU_B|.
    $$

    Úsečka $Q_1S_2$ prechádza postupne cez tri bunky susediace so stranou $AB$, teda cez $AP_1ES_2$, $P_1P_2FE$ a~$P_2BQ_1F$, takže na nej ležia body v~poradí $S_2$, $E$, $F$, $Q_1$. Keďže $U_A$ leží vnútri $S_2E$ a~$U_B$ vnútri $FQ_1$, platí
    $$
    d_{AB} = |U_AE| + |EF| + |FU_B|.
    $$

    Teraz každý z~týchto troch kúskov vyjadríme dotyčnicovými dĺžkami vo vrcholoch strednej bunky. Kružnica $w_a$ sa dotýka priamok $P_1R_2$ a~$Q_1S_2$, ktoré sa pretínajú v~bode $E$; nech $e$ je dĺžka dotyčnice z~$E$ ku kružnici $w_a$, takže $|U_AE| = e$. Analogicky nech $f$, $g$, $h$ sú dĺžky dotyčníc z~$F$, $G$, $H$ ku kružniciam $w_b$, $w_c$, $w_d$; vždy ide o~vrchol strednej bunky a~o~kružnicu tej rohovej bunky, ktorá tento vrchol obsahuje. Ďalej nech $e_0$, $f_0$, $g_0$, $h_0$ sú dĺžky dotyčníc z~$E$, $F$, $G$, $H$ k~$\omega$. Kružnica $\omega$ sa dotýka každej strany štvoruholníka $EFGH$ v~jej vnútornom bode, takže
    $$
    \gather
    |EF| = e_0 + f_0, \qquad |FG| = f_0 + g_0, \\
    |GH| = g_0 + h_0, \qquad |HE| = h_0 + e_0.
    \endgather
    $$
    Po dosadení je $d_{AB} = e + e_0 + f_0 + f$. Označme
    $$
    \gather
    x_E = e + e_0, \qquad x_F = f + f_0, \\
    x_G = g + g_0, \qquad x_H = h + h_0,
    \endgather
    $$
    čiže $d_{AB} = x_E + x_F$.

    Zvyšné tri strany vybavíme rovnako. Kružnice $w_b$, $w_c$ sa dotýkajú priamky $P_2R_1$ vo vnútorných bodoch úsečiek $P_2F$ a~$GR_1$, pričom ich bunky ležia v~tej polrovine určenej touto priamkou, ktorá obsahuje stranu $BC$; preto sú $BC$ aj $P_2R_1$ spoločné vonkajšie dotyčnice týchto dvoch kružníc a~$d_{BC}$ sa rovná vzdialenosti ich dotykových bodov na $P_2R_1$. Body na $P_2R_1$ idú v~poradí $P_2$, dotykový bod kružnice $w_b$, $F$, $G$, dotykový bod kružnice $w_c$, $R_1$, takže sa táto vzdialenosť rozkladá na dotyčnicu z~$F$ ku kružnici $w_b$, na $|FG|$ a~na dotyčnicu z~$G$ ku kružnici $w_c$. Dotyčnica z~$F$ ku kružnici $w_b$ má opäť dĺžku $f$, lebo $w_b$ sa dotýka oboch priamok prechádzajúcich bodom $F$, a~podobne pre $G$ a~$w_c$. Teda
    $$
    d_{BC} = f + f_0 + g_0 + g = x_F + x_G.
    $$
    Rovnako priamka $Q_2S_1$ dáva $d_{CD} = x_G + x_H$ a~priamka $P_1R_2$ dáva $d_{AD} = x_H + x_E$.

    Sčítaním dostávame
    $$
    \eqalign{
    d_{AB} + d_{CD} &= x_E + x_F + x_G + x_H = \cr
    &= d_{BC} + d_{AD}, \cr
    }
    $$
    čo je presne $(**)$. Tým je dokázané aj $(*)$ a~podľa Pitotovej vety má štvoruholník $ABCD$ vpísanú kružnicu.
}

\EnvId{EXLasGyo1ngIPJ1f2b0A2-kruznica-vpisana-do-uhla}
\Problem{0}{}{
    Nech $ABC$ je trojuholník s~polomerom vpísanej kružnice $r$ a~nech $\omega$ je kružnica s~polomerom $\rho < r$ vpísaná do uhla $BAC$. Druhé dotyčnice z~$B$ a~$C$ ku $\omega$ (iné než $AB$ a~$AC$) sa pretínajú v~bode $X$. Dokážte, že vpísaná kružnica trojuholníka $BCX$ sa dotýka vpísanej kružnice trojuholníka $ABC$.
}{
    Nech~$P$ je dotykový bod vpísanej kružnice trojuholníka $BCX$ s~$BC$. Chceme vyjadriť $|BP|$ pomocou $a$, $b$, $c$ (konkrétne dokázať $|BP|=s-b$).
}{
    Dotykové body vpísanej kružnice trojuholníka $BXC$ s~$BX$, $CX$ označme $K$, $L$ a~dotykové body $BX$ a~$CX$ s~$\omega$ ako $M$ a~$N$. Výpočet začína $|BP| = |BK| = |BM|-|KM|$.
}{
    Ďalší krok je uvedomiť si, že $|KM|=|LN|$ a~tiež $|BM|=|BW|$. To dáva $|BP| = |BW|-|LN|$. Dĺžka $|BW|$ je pekná, lebo je na strane $AB$. Ale $|LN|$ vieme tiež preniesť tam.
}{
    Všimni si $|LN|=|CN|-|CL| = |CV|-|CP|$. Takže máme $|BP|=|BW|-|CV|+|CP|$. Blížime sa.
}{
    Vezmi $|BW|=|AB|-|AW|$, $|CV|=|AC|-|AV|$, $|CP|=|BC|-|BP|$. Jeden jednoduchý krok od dokončenia.
}{
    Označme $\omega_0$ vpísanú kružnicu trojuholníka $ABC$ so stredom $I$, $O$ stred kružnice $\omega$, ďalej $W$, $V$ dotykové body $\omega$ s~priamkami $AB$, $AC$ a~$M$, $N$ dotykové body $\omega$ s~priamkami $BX$, $CX$. Vpísanú kružnicu trojuholníka $BCX$ označme $\omega_1$ a~jej dotykové body s~priamkami $BC$, $BX$, $CX$ postupne $P$, $K$, $L$. Celé riešenie stojí na rovnosti $|BP| = s - b$: vpísaná kružnica trojuholníka $ABC$ sa dotýka strany $BC$ vo vzdialenosti $s - b$ od $B$, takže z~nej vyplynie, že $\omega_0$ a~$\omega_1$ sa dotýkajú priamky $BC$ v~tom istom bode.

    Najprv si ujasnime polohu $\omega$. Kružnice $\omega$ a~$\omega_0$ sú obe vpísané do uhla $BAC$, takže ich stredy ležia na osi tohto uhla, a~ak $W_0$ je dotykový bod $\omega_0$ so stranou $AB$, sú pravouhlé trojuholníky $AWO$ a~$AW_0I$ podobné. Preto
    $$
    \frac{|AO|}{|AI|} = \frac{|AW|}{|AW_0|} = \frac{\rho}{r} < 1.
    $$
    Bod $O$ teda leží vnútri úsečky $AI$ a~z~rovnakých dotyčníc z~$A$ je
    $$
    |AW| = |AV| = \frac{\rho}{r}\,(s-a) < s-a.
    $$
    Z~trojuholníkovej nerovnosti je $s - a < b$ i~$s - a < c$, takže $W$ leží vnútri strany $AB$ a~$V$ vnútri strany $AC$, čiže $|BW| = c - |AW|$ a~$|CV| = b - |AV|$.

    Celý trojuholník $ABC$ leží v~páse medzi priamkou $BC$ a~rovnobežkou s~ňou vedenou bodom $A$, takže vzdialenosť jeho vnútorného bodu $I$ od $BC$, teda $r$, je menšia než vzdialenosť $A$ od $BC$. Vzdialenosť od priamky $BC$ sa pozdĺž úsečky $AI$ mení lineárne, a~preto je v~jej vnútornom bode $O$ väčšia než $r$, teda väčšia než $\rho$. Kružnica $\omega$ tak nemá s~priamkou $BC$ spoločný bod a~celá leží v~trojuholníku $ABC$; tu sme použili predpoklad $\rho < r$.

    Teraz konfigurácia. Priamka $BC$ nie je dotyčnicou $\omega$, takže $BX \neq BC$ a~$CX \neq BC$; zároveň $BX \neq AB$ a~$CX \neq AC$. Špeciálne bod $A$ neleží na $BX$ ani na $CX$, bod $C$ neleží na $BX$ a~bod $B$ neleží na $CX$. Bod $M$ leží na $\omega$, teda v~trojuholníku $ABC$, a~preto polpriamka $BM$ leží v~uhle $ABC$; nesplýva ani s~polpriamkou $BA$, ani s~polpriamkou $BC$, takže leží vo vnútri tohto uhla a~pretína úsečku $AC$ v~jej vnútornom bode. Priamka $BX$ preto oddeľuje $A$ od $C$ a~analogicky priamka $CX$ oddeľuje $A$ od $B$.

    Kružnica $\omega$ leží vo vnútri uhla, ktorý pri vrchole $B$ zvierajú jej dotyčnice, čiže polpriamky $BW$ a~$BM$, a~ten je celý v~polrovine s~hraničnou priamkou $BX$ obsahujúcej $A$; rovnako $\omega$ leží v~polrovine s~hraničnou priamkou $CX$ obsahujúcej $A$. Priamky $BX$ a~$CX$ delia rovinu na štyri oblasti a~podľa predošlého odseku je tá, ktorá obsahuje $A$, práve uhol vrcholovo protiľahlý k~uhlu $BXC$. Dotykový bod $M$ leží na priamke $BX$ aj v~uzávere tejto oblasti, čiže na polpriamke opačnej k~polpriamke $XB$, pričom $M \neq X$, lebo inak by dotyčnice $BX$ a~$CX$ splývali. Bod $X$ teda leží vnútri úsečky $BM$ a~rovnako vnútri úsečky $CN$. Špeciálne $X$ leží vo vnútri uhla $ABC$ i~uhla $ACB$, čiže vo vnútri trojuholníka $ABC$.

    Zvyšok je počítanie s~rovnakými dotyčnicami. Z~bodu $X$ máme $|XK| = |XL|$ (dotyčnice k~$\omega_1$) a~$|XM| = |XN|$ (dotyčnice k~$\omega$). Body $K$, $L$ ležia vnútri úsečiek $XB$, $XC$, kým $M$, $N$ ležia na opačných polpriamkach, takže
    $$
    |KM| = |XK| + |XM| = |XL| + |XN| = |LN|.
    $$
    Z~bodu $B$ je $|BM| = |BW|$ a~$|BK| = |BP|$, z~bodu $C$ je $|CN| = |CV|$ a~$|CL| = |CP|$. Bod $K$ leží vnútri úsečky $BM$ a~bod $L$ vnútri úsečky $CN$, preto
    $$
    \eqalign{
    |BP| &= |BK| = |BM| - |KM| = |BW| - |LN| \cr
         &= |BW| - |CN| + |CL| = |BW| - |CV| + |CP|. \cr
    }
    $$
    Dosadíme $|BW| = c - |AW|$, $|CV| = b - |AV|$ a~$|CP| = a - |BP|$, pričom $|AW| = |AV|$ sa vykráti:
    $$
    \eqalign{
    |BP| &= (c - |AW|) - (b - |AV|) + (a - |BP|) \cr
         &= a + c - b - |BP|. \cr
    }
    $$
    Odtiaľ $2|BP| = a + c - b$, čiže $|BP| = \frac{a+c-b}{2} = s - b$.

    Podľa vety o~dotykových bodoch vpísanej kružnice sa $\omega_0$ dotýka strany $BC$ v~bode $D$ s~$|BD| = s - b$. Body $P$ aj $D$ ležia vnútri úsečky $BC$ vo vzdialenosti $s - b$ od $B$, takže $P = D$ a~obe kružnice $\omega_0$, $\omega_1$ sa dotýkajú priamky $BC$ v~tom istom bode $P$. Ich stredy preto ležia na kolmici na $BC$ v~bode $P$, a~to na tej istej polpriamke: $\omega_0$ leží v~trojuholníku $ABC$ a~$\omega_1$ v~trojuholníku $BCX$, ktorý je vďaka tomu, že $X$ je vnútorný bod trojuholníka $ABC$, v~tej istej polrovine s~hraničnou priamkou $BC$. Kružnice sú rôzne, lebo $\omega_0$ sa dotýka $AB$ vo vnútornom bode $W_0$ tejto strany, ktorý neleží v~trojuholníku $BCX$; ten má totiž s~priamkou $AB$ spoločný jediný bod $B$. Ak označíme $r_1$ polomer $\omega_1$, je vzdialenosť stredov rovná $|r - r_1|$, pričom $r \neq r_1$, a~to je presne podmienka vnútorného dotyku. Kružnice $\omega_0$ a~$\omega_1$ sa teda dotýkajú, a~to v~bode $P$.
}

\EnvId{mZEAwZzSzf1cMcge-3Ovs-mnozina-priesecnikov-spolocnej-dotycnice}
\Problem{0}{USAMO 1991}{
    Nech $D$ je bod na strane $BC$ trojuholníka $ABC$. Do trojuholníkov $BDA$ a~$DCA$ vpíšme kružnice. Ich spoločná vonkajšia dotyčnica iná než $BC$ pretína $AD$ v~bode $X$. Určte množinu bodov $X$, keď $D$ prebieha po strane $BC$.
}{
    Kľúčom je správna hypotéza -- uváž hraničné prípady, keď $D=B$ a~$D=C$. Nenaznačuje to niečo?
}{
    V špeciálnych prípadoch vychádza $|AX|$ rovnako, konkrétne rovné $s-a$. Možno je odpoveďou kružnica so stredom $A$ a~týmto polomerom?
}{
    Dotykové body vpísanej kružnice trojuholníka $ABD$ s~$BD$, $AB$ označme $P$, $R$; dotykové body vpísanej kružnice trojuholníka $ACD$ s~$CD$, $CA$ označme $Q$, $S$. Vieme, že chceme dokázať $|AX|=s-a$. Prvé zjednodušenie je použiť $|AX|=|AD|-|DX|$ a~spomenúť si na jednu z~viet, ktorá hovorí, čomu sa $|DX|$ rovná.
}{
    Máme $|DX|=|PQ|$, takže $|AX|=|AD|-|DX|=|AD|-|PQ|$. S~$|AD|$ už viac nespravíme, tak sa zamerajme na $|PQ|$.
}{
    Máme $|PQ| = |DP|+|DQ|$, a~teda $|AX| = |AD|-|DP|-|DQ|$. Ale to je v~poriadku. Čomu sa rovnajú $|PD|$ a~$|DQ|$?
}{
    Dĺžku $|DP|$ vieme vyjadriť pomocou vzorca typu $s-a$ v~$ABD$, podobne $|QD|$.
}{
    \Answer{Hľadanou množinou je otvorený oblúk kružnice so stredom $A$ a~polomerom $s - a = \frac{b + c - a}{2}$ ležiaci vnútri uhla $BAC$. Jeho krajnými bodmi sú dotykové body vpísanej kružnice trojuholníka $ABC$ so stranami $AB$ a~$AC$, tie však do hľadanej množiny nepatria.}
    Odpoveď uhádneme z~hraničných prípadov. Keď sa $D$ blíži k~$B$, vpísaná kružnica trojuholníka $ABD$ sa zmršťuje do bodu $B$ a~vpísaná kružnica trojuholníka $ADC$ prechádza na vpísanú kružnicu trojuholníka $ABC$. Spoločnými vonkajšími dotyčnicami sa tak v~limite stávajú dve dotyčnice z~bodu $B$ ku vpísanej kružnici trojuholníka $ABC$, čiže priamky $BC$ a~$AB$, a~bod $X$ dobehne do dotykového bodu tejto kružnice so stranou $AB$. Ten je od $A$ vzdialený $s - a$. Symetricky pre $D \to C$ vyjde dotykový bod na $AC$, opäť vo vzdialenosti $s - a$ od $A$. Hypotéza teda znie $|AX| = s - a$ pre každú polohu $D$.

    Nech je ďalej $D$ vnútorným bodom strany $BC$; pre $D = B$ alebo $D = C$ jeden z~trojuholníkov zo zadania degeneruje. Vpísanú kružnicu trojuholníka $ABD$ označme $k_1$, vpísanú kružnicu trojuholníka $ADC$ označme $k_2$ a~ich dotykové body s~priamkou $BC$ postupne $P$ a~$Q$. Obe kružnice sa dotýkajú priamky $BC$ a~ležia v~tej istej polrovine, takže $BC$ je ich spoločná vonkajšia dotyčnica; druhú spoločnú vonkajšiu dotyčnicu označme $q$. Obe sa dotýkajú aj priamky $AD$, pričom ležia v~opačných polrovinách určených touto priamkou, takže $AD$ je ich spoločná vnútorná dotyčnica (a~kružnice sú disjunktné, prípadne sa dotýkajú v~bode ležiacom na $AD$; v~tom prípade je $AD$ ich jedinou spoločnou vnútornou dotyčnicou a~nasledujúca veta, ktorej dôkaz pracuje iba s~rovnakými dotyčnicami z~$D$ a~z~$X$, platí bez zmeny). Priamka $AD$ pretína $BC$ v~bode $D$ a~priamku $q$ v~bode $X$, takže podľa vety o~spoločnej vnútornej dotyčnici dvoch kružníc je
    $$
    |DX| = |PQ|.
    $$

    Bod $P$ leží vnútri úsečky $BD$ a~bod $Q$ vnútri úsečky $DC$, preto $|PQ| = |DP| + |DQ|$. Vzorce typu $s - a$ použité pri vrchole $D$ v~trojuholníku $ABD$, resp. $ADC$ (protiľahlé strany sú $AB$, resp. $AC$) dávajú
    $$
    |DP| = \frac{|AD| + |BD| - c}{2}, \qquad |DQ| = \frac{|AD| + |CD| - b}{2}.
    $$
    Keďže $|BD| + |CD| = a$, ich sčítaním dostaneme
    $$
    |DX| = |PQ| = |AD| - \frac{b + c - a}{2} = |AD| - (s - a).
    $$

    Ostáva overiť polohu bodu $X$. Kružnice $k_1$, $k_2$ ležia v~polrovine ohraničenej priamkou $BC$ a~obsahujúcej $A$, takže v~nej ležia aj ich dotykové body s~priamkou $q$. Tie sú po opačných stranách priamky $AD$: každý leží na svojej kružnici a~na priamke $AD$ ležať nemôže, lebo inak by sa $q$ dotýkala tejto kružnice v~jej dotykovom bode s~$AD$, čiže by bolo $q = AD$, čo je vylúčené (priamka $AD$ obe kružnice oddeľuje, $q$ nie). Bod $X$ ako priesečník priamky $q$ s~priamkou $AD$ preto leží medzi nimi, a~teda v~spomínanej polrovine, čiže na polpriamke $DA$. Spolu s~$0 < |DX| < |AD|$ to znamená, že $X$ je vnútorným bodom úsečky $AD$, a~tak
    $$
    |AX| = |AD| - |DX| = s - a.
    $$

    Bod $X$ teda vždy leží na kružnici $\omega$ so stredom $A$ a~polomerom $s - a$, a~keďže je vnútorným bodom úsečky $AD$ pre vnútorný bod $D$ strany $BC$, leží vnútri uhla $BAC$. Naopak, každý bod tohto oblúka aj nastane: keď $D$ prebieha vnútro strany $BC$, polpriamka $AD$ prebehne práve všetky polpriamky z~$A$ ležiace ostro medzi polpriamkami $AB$ a~$AC$, a~to vzájomne jednoznačne, pričom na každej z~nich je jediný bod vo vzdialenosti $s - a$ od $A$ a~podľa dokázaného je ním práve bod $X$ prislúchajúci tomuto $D$. Krajnými bodmi oblúka sú dotykové body vpísanej kružnice trojuholníka $ABC$ so stranami $AB$ a~$AC$, lebo tie sú od $A$ vzdialené presne $s - a$; zodpovedajú degenerovaným polohám $D = B$ a~$D = C$, a~preto do hľadanej množiny nepatria.
}

\EnvId{WqLo9VdtCR-oLjiZTe-hJ-dotycnicovy-stvoruholnik-a-opisana-kruznica}
\Problem{1}{IMO 2021, P4}{
    Nech $ABCD$ je dotyčnicový štvoruholník s~vpísanou kružnicou so stredom $I$ a~nech $\Omega$ je opísaná kružnica trojuholníka $AIC$. Predĺženia $BA$ a~$BC$ za body $A$ a~$C$ pretínajú $\Omega$ v~bodoch $X$ a~$Z$ a~predĺženia $AD$ a~$CD$ za bodom $D$ pretínajú $\Omega$ v~bodoch $Y$ a~$T$. Dokážte, že
    $$
    |AD| + |DT| + |TX| + |XA| = |CD| + |DY| + |YZ| + |ZC|.
    $$
}{
    Máme dotyčnice, tak sa zdá, že vieme urobiť trochu mágie s~rovnakými dotyčnicami. Skoro, treba si niečo uvedomiť. Skús naslepo odvodzovať pekné vlastnosti.
}{
    Kľúčom je nájsť zhodné trojuholníky. Použi kružnicu prechádzajúcu bodmi $A, I, C, X, Y, T, Z$, aby si na to prišiel.
}{
    Nech $P, Q, R, S$ sú dotykové body vpísanej kružnice na $AB$, $BC$, $CD$, $DA$. Skús dokázať, že trojuholníky $IQZ$ a~$IRT$ sú zhodné.
}{
    Rovnosť uhlov pri $Z$ a~$T$ vieme dostať počítaním uhlov. Uhly pri $Q$ a~$R$ sú pravé, takže zhodnosť nasleduje. Existuje niekde inde podobná zhodnosť?
}{
    V skutočnosti aj $IPX$ a~$ISY$ sú zhodné. Tieto dve zhodnosti spolu dávajú rovnosti úsečiek potrebné na záverečný výpočet.
}{
    Označme $r$ polomer vpísanej kružnice štvoruholníka $ABCD$, ďalej $P$, $Q$, $R$, $S$ jej dotykové body postupne na stranách $AB$, $BC$, $CD$, $DA$ a~$O$ stred kružnice $\Omega$. Bod $I$ má od každej z~priamok $AB$, $BC$, $CD$, $DA$ vzdialenosť $r > 0$, takže na žiadnej z~nich neleží; špeciálne je rôzny od bodov $X$, $Y$, $Z$, $T$.

    Dotykové body sú vnútornými bodmi príslušných strán a~podľa zadania ležia $X$, $Z$ za bodmi $A$, resp. $C$ a~body $Y$, $T$ za bodom $D$. Na priamke $AB$ tak ležia body v~poradí $B$, $P$, $A$, $X$, na priamke $BC$ v~poradí $B$, $Q$, $C$, $Z$, na priamke $AD$ v~poradí $A$, $S$, $D$, $Y$ a~na priamke $CD$ v~poradí $C$, $R$, $D$, $T$. O~tieto štyri poradia sa opiera každý ďalší krok.

    Samotné rovnaké dotyčnice nestačia: väčšina úsečiek zo zadania vôbec neleží na stranách štvoruholníka. Mostom k~nim sú dve dvojice zhodných pravouhlých trojuholníkov s~vrcholom $I$, v~ktorých pravé uhly dodajú dotykové body vpísanej kružnice a~zhodné ostré uhly obvodové uhly v~$\Omega$.

    \textbf{Trojuholníky $IQZ$ a~$IRT$.} Polomer do dotykového bodu je kolmý na dotyčnicu, takže $IQ \perp BC$ a~$IR \perp CD$. Bod $Z$ leží na priamke $BC$ a~$Z \neq Q$, bod $T$ na priamke $CD$ a~$T \neq R$, takže $IQZ$ a~$IRT$ sú pravouhlé trojuholníky s~pravými uhlami pri $Q$, resp. pri $R$, a~s~odvesnami $|IQ| = |IR| = r$.

    Z~poradia $B$, $Q$, $C$, $Z$ ležia body $Q$ a~$C$ na tej istej polpriamke z~bodu $Z$, takže $|\angle IZQ| = |\angle IZC|$; z~poradia $C$, $R$, $D$, $T$ ležia body $R$ a~$C$ na tej istej polpriamke z~bodu $T$, takže $|\angle ITR| = |\angle ITC|$. Oba uhly sú ostré, lebo sú to nepravé uhly pravouhlých trojuholníkov. Body $I$, $C$, $Z$, $T$ ležia na $\Omega$, takže $|\angle IZC|$ a~$|\angle ITC|$ sú obvodové uhly nad tetivou $IC$: buď sú zhodné (keď $Z$ a~$T$ ležia na tom istom oblúku určenom bodmi $I$, $C$), alebo je ich súčet $180^\circ$. Druhá možnosť by však znamenala, že aspoň jeden z~nich nie je ostrý, čiže
    $$
    |\angle IZQ| = |\angle IZC| = |\angle ITC| = |\angle ITR|.
    $$

    Odtiaľ $|\angle QIZ| = 90^\circ - |\angle IZQ| = 90^\circ - |\angle ITR| = |\angle RIT|$. Trojuholníky $IQZ$ a~$IRT$ majú teda zhodnú stranu $|IQ| = |IR|$ a~zhodné uhly pri oboch jej krajných bodoch, takže sú zhodné, čo dáva
    $$
    |QZ| = |RT|, \qquad |IZ| = |IT|.
    $$

    \textbf{Trojuholníky $IPX$ a~$ISY$.} Rovnako je $IP \perp AB$, $IS \perp DA$ a~$|IP| = |IS| = r$, pričom $X \neq P$ a~$Y \neq S$, takže $IPX$ a~$ISY$ sú pravouhlé trojuholníky s~pravými uhlami pri $P$, resp. pri $S$. Z~poradia $B$, $P$, $A$, $X$ je $|\angle IXP| = |\angle IXA|$ a~z~poradia $A$, $S$, $D$, $Y$ je $|\angle IYS| = |\angle IYA|$. Oba tieto uhly sú opäť ostré a~sú to obvodové uhly nad tetivou $IA$ kružnice $\Omega$, takže sa rovnajú. Tým istým argumentom ako vyššie sú trojuholníky $IPX$ a~$ISY$ zhodné a
    $$
    |PX| = |SY|, \qquad |IX| = |IY|.
    $$

    \textbf{Prenos na tetivy $TX$ a~$YZ$.} Bod $I$ leží na $\Omega$, takže $|OI|$ je polomer $\Omega$ a~$O \neq I$. Ďalej $X \neq Y$: inak by tento bod ležal na priamkach $AB$ aj $AD$, ktoré majú jediný spoločný bod $A$, no $X \neq A$. Body $O$ aj $I$ sú rovnako vzdialené od $X$ ako od $Y$, takže priamka $OI$ je osou úsečky $XY$. Podobne $Z \neq T$ (inak by tento bod ležal na priamkach $BC$ aj $CD$, teda by to bol bod $C$, pričom $T \neq C$) a~z~$|OZ| = |OT|$, $|IZ| = |IT|$ je priamka $OI$ aj osou úsečky $ZT$. Osová súmernosť podľa priamky $OI$ preto zobrazí $X$ na $Y$ a~$T$ na $Z$, odkiaľ
    $$
    |TX| = |YZ|.
    $$

    \textbf{Záverečný výpočet.} Podľa vety o~dvoch dotyčniciach z~bodu je $|AP| = |AS|$, $|CQ| = |CR|$ a~$|DR| = |DS|$. Zo štyroch poradí bodov ďalej dostávame
    $$
    \gather
    |AD| = |AS| + |SD|, \qquad |CD| = |CR| + |RD|, \\
    |XA| = |PX| - |PA|, \qquad |ZC| = |QZ| - |QC|, \\
    |DT| = |RT| - |RD|, \qquad |DY| = |SY| - |SD|.
    \endgather
    $$
    Po dosadení sa vďaka $|AS| = |PA|$, resp. $|CR| = |QC|$ vykrátia dotyčnicové úseky pri $A$, resp. pri $C$ a~zostane
    $$
    \eqalign{
        |AD| + |DT| + |XA| &= |SD| - |RD| + |RT| + |PX|,\cr
        |CD| + |DY| + |ZC| &= |RD| - |SD| + |QZ| + |SY|.\cr
    }
    $$
    Pravé strany sú rovnaké, lebo $|SD| = |RD|$, $|RT| = |QZ|$ a~$|PX| = |SY|$. Pripočítaním rovnosti $|TX| = |YZ|$ dostávame
    $$
    |AD| + |DT| + |TX| + |XA| = |CD| + |DY| + |YZ| + |ZC|.
    $$
}

\EnvId{zwCIF1ytYOGA5L7Ai1rsb-obrazy-dotykovych-bodov-podla-stredu}
\Problem{1}{IMO Shortlist 2005, G1}{
    Nech $ABC$ je trojuholník so stredom vpísanej kružnice $I$, ktorý spĺňa $|AC| + |BC| = 3|AB|$, a~nech $D$, $E$ sú body, v~ktorých sa vpísaná kružnica dotýka strán $BC$ a~$CA$. Nech $K$ a~$L$ sú obrazy bodov $D$ a~$E$ v~stredovej súmernosti podľa $I$. Dokážte, že body $A$, $B$, $K$, $L$ ležia na jednej kružnici.
}{
    Prvú nápovedu dostaneme po tom, ako nakreslíme dokonalý obrázok a~uvedomíme si, že $I$ leží na kružnici s~$A, B, K, L$. A~kružnica cez $AIB$ nám dá nápovedu, čo nakresliť ďalej: čím ešte prechádza?
}{
    Kružnica $A, I, B$ prechádza stredom kružnice pripísanej oproti $C$, teda $I_c$. Nakreslime teda kružnicu pripísanú oproti $C$.
}{
    Záhadná podmienka $|AC|+|BC|=3|AB|$ dáva väčší zmysel, keď zavedieme dotykové body kružnice pripísanej oproti $C$ s~$CB$ a~$CA$. Čas na trochu počítania s~rovnakými dotyčnicami.
}{
    Ukáž, že naša podmienka je ekvivalentná tomu, že $D$ je stredom $CD'$, podobne pre $E$. Aké sú dôsledky v~obrázku?
}{
    Zdôvodni, že $|I_cD'| = 2|ID|$. To konečne dáva nápovedu, čo robiť s~$K$. Teraz sme blízko k~dokončeniu.
}{
    Presný obrázok napovedá, že na hľadanej kružnici leží aj bod $I$. To je dobrá stopa: kružnica opísaná trojuholníku $AIB$ je známa konfigurácia, prechádza totiž stredom $I_c$ kružnice pripísanej oproti vrcholu $C$ a~úsečka $II_c$ je jej priemerom. Dokážeme preto, že $A$, $B$, $K$, $L$ ležia na kružnici $\omega$ s~priemerom $II_c$.

    Bod $I_c$ leží na osi vnútorného uhla pri vrchole $C$ a~na osiach vonkajších uhlov pri vrcholoch $A$ a~$B$. Priamky $AI$ a~$AI_c$ sú teda os vnútorného a~os vonkajšieho uhla pri $A$, ktoré sú na seba kolmé, takže $|\angle IAI_c| = 90^\circ$, a~rovnako $|\angle IBI_c| = 90^\circ$. Podľa Tálesovej vety ležia $A$ i~$B$ na $\omega$.

    Nech sa kružnica pripísaná oproti $C$ dotýka priamok $CB$ a~$CA$ v~bodoch $D'$ a~$E'$. Táto kružnica sa dotýka strany $AB$ a~predĺžení strán $CB$ a~$CA$ za bodmi $B$ a~$A$, takže $D'$ leží na polpriamke $CB$ a~$E'$ na polpriamke $CA$; na tých istých polpriamkach ležia aj dotykové body $D$, $E$ vpísanej kružnice. Vzdialenosti všetkých štyroch od vrcholu $C$ poznáme:
    $$
    |CD'| = |CE'| = s, \qquad |CD| = |CE| = s - c.
    $$
    Podmienka $|AC| + |BC| = 3|AB|$ hovorí $b + a = 3c$, teda $a + b + c = 4c$, čiže $s = 2c$. To je presne rovnosť $s = 2(s - c)$, a~teda
    $$
    |CD'| = 2|CD|, \qquad |CE'| = 2|CE|.
    $$
    Keďže $D$ aj $D'$ ležia na polpriamke $CB$ a~$E$ aj $E'$ na polpriamke $CA$, znamená to, že $D$ je stredom úsečky $CD'$ a~$E$ je stredom úsečky $CE'$.

    Body $I$ a~$I_c$ ležia obidva na osi uhla $ACB$, čiže na polpriamke $CI$, a~úsečky $ID$, $I_cD'$ sú kolmé na $CB$. Pravouhlé trojuholníky $CDI$ a~$CD'I_c$ sa preto zhodujú v~uhle pri vrchole $C$ a~sú podobné s~koeficientom $|CD'| / |CD| = 2$, odkiaľ
    $$
    |I_cD'| = 2|ID|, \qquad |CI_c| = 2|CI|.
    $$
    Body $I$, $I_c$ ležia na tej istej polpriamke z~$C$, takže druhá rovnosť hovorí, že $I$ je stredom úsečky $CI_c$; zvlášť $I \neq I_c$ a~kružnica $\omega$ je naozaj kružnica.

    Stredová súmernosť podľa $I$ zobrazuje $D$ na $K$, $E$ na $L$ a~podľa predchádzajúceho aj $C$ na $I_c$. Úsečka $KI_c$ je teda obrazom úsečky $DC$, čiže
    $$
    KI_c \parallel CB, \qquad |KI_c| = |CD| = s - c > 0,
    $$
    a~rovnako $LI_c \parallel CA$ s~$|LI_c| = |CE| > 0$. Zároveň je $I$ stredom úsečky $DK$, takže priamka $IK$ je priamkou $DK$, ktorá je kolmá na $CB$; pritom $|IK| = |ID| > 0$. V~bode $K$ sa tak stretá kolmica na $CB$ s~rovnobežkou s~$CB$, čo dáva $|\angle IKI_c| = 90^\circ$; rovnako z~$IL \perp CA$ a~$LI_c \parallel CA$ dostaneme $|\angle ILI_c| = 90^\circ$.

    Podľa Tálesovej vety ležia teda aj body $K$ a~$L$ na kružnici $\omega$ s~priemerom $II_c$, na ktorej už ležia $A$ a~$B$.
}

\EnvId{cqRxrMgyFJE7-aHA9USqB-rozdelenie-obvodu-trojuholnika}
\Problem{2}{}{
    Nech $ABC$ je trojuholník s~obvodom $4$. Body $X$ a~$Y$ ležia na polpriamkach $AB$ a~$AC$ tak, že $|AX| = |AY| = 1$ a~úsečky $BC$ a~$XY$ sa pretínajú v~bode $F$. Dokážte, že jeden z~dvoch trojuholníkov, na ktoré $AF$ rozdelí trojuholník $ABC$, má obvod $2$.
}{
    Prvý trik je spomenúť si na jedno obzvlášť pekné miesto, kde sa v~geometrii trojuholníka nejako objavuje obvod. Určite sme ho v~tomto letáku videli.
}{
    Vzdialenosť od $A$ k~dotykovým bodom kružnice pripísanej oproti $A$ na priamkach $AB$ a~$AC$ je presne polovica obvodu, teda 2. A~$|AX|=|AY|=1$. Je to náhoda?
}{
    Nech sa kružnica pripísaná oproti $A$ dotýka $AB$ a~$AC$ v~$D$ a~$E$. Máme $|AD|=|AE|=2$ a~$|AX|=|AY|=1$, a~tak je priamka $XY$ obrazom $A$-poláry kružnice pripísanej oproti $A$ v~rovnoľahlosti so stredom $A$ a~koeficientom $1/2$. Čo o~tejto priamke vieme?
}{
    Priamka $XY$ je chordála bodu $A$ a~kružnice pripísanej oproti $A$! A~$F$ na nej leží. Blížime sa.
}{
    Posledný bod, ktorý nakreslíme, je dotykový bod pripísanej kružnice a~$BC$. Použitím faktu, že $F$ leží na tejto chordále, dostaneme nejaké pekné rovnaké dĺžky.
}{
    Obvod trojuholníka $ABC$ je $4$, takže $s = 2$. Polovica obvodu sa v~trojuholníku objavuje ako vzdialenosť od $A$ k~dotykovým bodom kružnice pripísanej oproti vrcholu $A$, a~hodnota $|AX| = |AY| = 1$ je presne jej polovica. To hovorí, čo dokresliť.

    Označme $\omega$ kružnicu pripísanú trojuholníku $ABC$ oproti vrcholu $A$, jej stred $J$ a~jej polomer $\rho$, a~nech sa $\omega$ dotýka priamok $AB$, $AC$, $BC$ postupne v~bodoch $D$, $E$, $T$. Podľa vety o~pripísanej kružnici je $|AD| = |AE| = s = 2$. Body $D$, $X$ ležia na polpriamke $AB$ a~body $E$, $Y$ na polpriamke $AC$, takže z~$|AX| = |AY| = 1$ je $X$ stredom úsečky $AD$ a~$Y$ stredom úsečky $AE$. Priamka $XY$ je teda obrazom priamky $DE$ v~rovnoľahlosti $h$ so stredom $A$ a~koeficientom $\frac12$.

    Kľúčové tvrdenie je, že pre každý bod $P$ priamky $XY$ platí $|PA|^2 = |PJ|^2 - \rho^2$, teda že $XY$ je chordálou bodu $A$ (chápaného ako kružnica s~nulovým polomerom) a~kružnice $\omega$.

    Dokážeme to. Body $D$, $E$ sú rôzne, lebo inak by ležali na priamke $AB$ aj na priamke $AC$, a~teda by splynuli s~$A$. Z~$|AD| = |AE|$ a~$|JD| = |JE| = \rho$ ležia $A$ aj $J$ na osi úsečky $DE$, takže $AJ \perp DE$. Nech $H$ je priesečník priamok $DE$ a~$AJ$ a~nech $G$ je stred úsečky $AH$. Rovnoľahlosť $h$ zobrazuje $H$ na $G$, takže $G$ leží na $XY$; priamka $XY$ je rovnobežná s~$DE$, teda kolmá na $AJ$, a~pretína ju práve v~bode $G$.

    Dotyčnica $AD$ je kolmá na polomer $JD$, takže trojuholník $ADJ$ má pri $D$ pravý uhol a~$DH$ je jeho výška na preponu $AJ$. Pytagorova veta dáva $|AJ|^2 = |AD|^2 + \rho^2$ a~Euklidova veta o~odvesne dáva
    $$
    |AD|^2 = |AH| \cdot |AJ| = 2|AG| \cdot |AJ|.
    $$
    Päta výšky $H$ leží vnútri prepony $AJ$, takže tam leží aj $G$ a~platí $|GJ| = |AJ| - |AG|$. Pre bod $P$ na priamke $XY$ je $|PA|^2 = |PG|^2 + |GA|^2$ a~$|PJ|^2 = |PG|^2 + |GJ|^2$ (Pytagorova veta, pre $P = G$ triviálne), a~teda
    $$
    \eqalign{
    |PJ|^2 - |PA|^2 &= |GJ|^2 - |GA|^2 = \cr
    &= \bigl(|AJ| - |AG|\bigr)^2 - |AG|^2 = \cr
    &= |AJ|^2 - 2|AG| \cdot |AJ| = |AJ|^2 - |AD|^2 = \rho^2, \cr
    }
    $$
    čo je práve dokazované tvrdenie.

    Bod $F$ leží na priamke $XY$, takže $|FA|^2 = |FJ|^2 - \rho^2$. Zároveň leží na priamke $BC$, ktorá je dotyčnicou $\omega$ v~bode $T$, takže $JT \perp BC$ a~z~Pytagorovej vety je $|FJ|^2 = |FT|^2 + \rho^2$. Spolu
    $$
    |FA| = |FT|.
    $$

    Podľa vety o~dotykových bodoch pripísanej kružnice pri zvyšných dvoch vrcholoch je $|BT| = s - c$ a~$|CT| = s - b$. Obe dĺžky sú kladné a~ich súčet je $2s - b - c = a = |BC|$, takže $T$ leží vnútri úsečky $BC$. Bod $F$ leží na úsečke $BC$ a~je jej vnútorným bodom, inak by $AF$ trojuholník $ABC$ nerozdelila na dva trojuholníky. Nastáva teda jeden z~dvoch prípadov.

    \begitems
    \i Bod $F$ leží na úsečke $BT$. Potom $|BF| + |FT| = |BT|$ a~obvod trojuholníka $ABF$ je
    $$
    \eqalign{
    |AB| + |BF| + |FA| &= c + |BF| + |FT| = \cr
    &= c + |BT| = c + (s - c) = 2. \cr
    }
    $$
    \i Bod $F$ leží na úsečke $TC$. Potom $|CF| + |FT| = |CT|$ a~obvod trojuholníka $ACF$ je
    $$
    \eqalign{
    |AC| + |CF| + |FA| &= b + |CF| + |FT| = \cr
    &= b + |CT| = b + (s - b) = 2. \cr
    }
    $$
    \enditems
}

\EnvId{p_QDQdc3MiS1qjG5JIegA-brianchonova-veta}
\Problem{2}{Brianchonova veta}{
    Ak je šesťuholník $ABCDEF$ opísaný kružnici (t.j. každá z~jeho šiestich strán sa dotýka kružnice), tak jeho tri hlavné uhlopriečky $AD$, $BE$, $CF$ sa pretínajú v~jednom bode.
}{
    Nakresli tri kružnice tak, aby uhlopriečky $AD$, $BE$, $CF$ boli chordály. Je to zákerné.
}{
    Jedna z~kružníc sa dotýka polpriamok opačných k~$BA$ a~$DE$. Analogicky ostatné. Jediný trik je umiestniť ich tak, aby chordály fungovali.
}{
    Skús ich umiestniť tak, aby dĺžky vonkajších dotyčníc ku každej kružnici a~vpísanej kružnici $ABCDEF$ boli rovnaké.
}{
    Tri priamky prechádzajúce jedným bodom volajú po potenčnom strede: hľadáme teda tri kružnice, ktorých chordálami sú práve uhlopriečky $AD$, $BE$, $CF$. Mocnosť bodu $M$ ku kružnici $\omega'$ značme $p(M,\omega')$ a~pripomeňme, že pre $M$ na dotyčnici kružnice $\omega'$ s~dotykovým bodom $Z$ je $p(M,\omega') = |MZ|^2$. Mocnosti vrcholov teda vieme uchopiť cez dotyčnice, a~tie sú to jediné, čo v~úlohe máme.

    Nech sa vpísaná kružnica $\omega$ so stredom $O$ a~polomerom $r$ dotýka strán $AB$, $BC$, $CD$, $DE$, $EF$, $FA$ postupne v~bodoch $P$, $Q$, $R$, $S$, $T$, $U$. V~každom vrchole zvierajú obe tamojšie strany uhol, vnútri ktorého leží $\omega$, takže všetky vnútorné uhly sú menšie než $180^\circ$ a~šesťuholník je konvexný; body $P$, $Q$, $R$, $S$, $T$, $U$ preto ležia na $\omega$ v~tomto poradí, v~rovnakom smere obehu ako vrcholy. Podľa Tvrdenia 1 sú obe dotyčnice z~každého vrcholu rovnako dlhé, položme
    $$
    \gather
    a = |AP| = |AU|, \quad b = |BP| = |BQ|, \quad c = |CQ| = |CR|,\\
    d = |DR| = |DS|, \quad e = |ES| = |ET|, \quad f = |FT| = |FU|.
    \endgather
    $$

    Hľadané kružnice budú dotyčnicami priamok protiľahlých strán: nech sa $\omega_1$ dotýka priamok $AB$, $DE$ v~bodoch $P_1$, $S_1$, kružnica $\omega_2$ priamok $BC$, $EF$ v~bodoch $Q_1$, $T_1$ a~kružnica $\omega_3$ priamok $CD$, $FA$ v~bodoch $R_1$, $U_1$. Každý vrchol potom leží na dotyčnici práve dvoch z~nich, napríklad $B$ na dotyčniciach $AB$, $BC$ kružníc $\omega_1$, $\omega_2$, a~priamka $BE$ bude ich chordálou, len čo $|BP_1| = |BQ_1|$ a~$|ES_1| = |ET_1|$. Chceme teda, aby v~každom vrchole boli oba tamojšie nové dotykové body od neho rovnako vzdialené.

    Pôvodné dotykové body to spĺňajú, lenže dávajú $\omega_1 = \omega_2 = \omega_3 = \omega$. Posunieme ich preto po stranách o~tú istú dĺžku, striedavo v~smere obehu a~proti nemu: na každej strane leží práve jeden z~vrcholov $B$, $D$, $F$ a~dotykový bod posunieme smerom k~nemu, a~to až zaň. Zvoľme $t > \max(b,d,f)$ a~nech body $P_1$, $Q_1$ ležia na polpriamkach opačných k~$BA$, $BC$ vo vzdialenosti $t-b$ od $B$, body $R_1$, $S_1$ na polpriamkach opačných k~$DC$, $DE$ vo vzdialenosti $t-d$ od $D$ a~body $T_1$, $U_1$ na polpriamkach opačných k~$FE$, $FA$ vo vzdialenosti $t-f$ od $F$. Potom
    $$
    \gather
    |AP_1| = |AU_1| = a+t, \quad |BP_1| = |BQ_1| = t-b,\\
    |CQ_1| = |CR_1| = c+t, \quad |DR_1| = |DS_1| = t-d,\\
    |ES_1| = |ET_1| = e+t, \quad |FT_1| = |FU_1| = t-f,
    \endgather
    $$
    napríklad $|AP_1| = |AB| + |BP_1| = (a+b)+(t-b)$ a~$|CQ_1| = |CB| + |BQ_1| = (b+c)+(t-b)$. Zároveň každý nový bod vznikol posunutím toho pôvodného po jeho strane o~dĺžku práve $t$, napríklad $|PP_1| = |PB| + |BP_1| = t$.

    Ostáva ukázať, že kružnica dotýkajúca sa $AB$ v~bode $P_1$ a~$DE$ v~bode $S_1$ vôbec existuje; práve tu sa striedanie smerov spotrebuje. Nech $m_1$ je os úsečky $PS$ a~$\sigma$ osová súmernosť podľa nej. Z~$|OP| = |OS| = r$ prechádza $m_1$ bodom $O$, takže $\sigma$ zobrazuje $\omega$ na seba, bod $P$ na bod $S$, a~teda dotyčnicu $AB$ na dotyčnicu $DE$. Ako nepriama zhodnosť obracia smer obehu: smer $A \to B$ je v~bode $P$ smerom obehu dopredu, jeho obrazom je preto v~bode $S$ smer obehu dozadu, čiže smer $E \to D$. Body $P_1$, $S_1$ pritom vznikli posunom o~$t$ práve v~týchto dvoch smeroch, takže $\sigma(P_1) = S_1$.

    Priamka $m_1$ nie je kolmá na $AB$: inak by splynula s~priamkou $OP$, jedinou kolmicou na $AB$ vedenou bodom $O$, a~$\sigma$ by nechala $P$ na mieste, čiže $P = S$. Kolmica na $AB$ v~bode $P_1$ preto pretne $m_1$ v~jedinom bode $O_1$; ten $\sigma$ nechá na mieste a~túto kolmicu zobrazí na kolmicu na $DE$ v~bode $S_1$, takže $|O_1P_1| = |O_1S_1|$ a~kružnica $\omega_1$ so stredom $O_1$ a~polomerom $|O_1P_1|$ je hľadaná. Nulový polomer by znamenal $O_1 = P_1 = S_1$, teda že $P_1$ je priesečníkom priamok $AB$ a~$DE$ (rôzne sú, lebo sa $\omega$ dotýkajú v~rôznych bodoch). Bod $P_1$ však leží na $AB$ vo vzdialenosti $t$ od $P$ v~pevnom smere, takže s~týmto jediným priesečníkom splynie najviac pre jednu hodnotu $t$; pri troch dvojiciach protiľahlých strán to vylúči najviac tri hodnoty a~$t$ volíme mimo nich. Rovnako zostrojíme $\omega_2$ so stredom $O_2$ na osi $m_2$ úsečky $QT$ a~$\omega_3$ so stredom $O_3$ na osi $m_3$ úsečky $RU$.

    Stredy sú navzájom rôzne. Bod $O$ leží na kolmici na $AB$ v~bode $P$, kým $O_1$ na rovnobežnej kolmici v~bode $P_1 \neq P$, takže $O_1 \neq O$, a~rovnako $O_2 \neq O$, $O_3 \neq O$. Osi $m_1$, $m_2$, $m_3$ pritom prechádzajú bodom $O$ a~sú navzájom rôzne: body $P$, $Q$, $S$, $T$ ležia na $\omega$ v~tomto cyklickom poradí, takže sa tetivy $PS$ a~$QT$ pretínajú, nie sú teda rovnobežné a~nesplývajú ani kolmice na ne vedené bodom $O$; analogicky pre dvojice $QT$, $RU$ a~$RU$, $PS$. Priamky $m_1$, $m_2$, $m_3$ tak majú jediný spoločný bod $O$, a~keďže $O_i$ leží na $m_i$ a~$O_i \neq O$, sú stredy $O_1$, $O_2$, $O_3$ navzájom rôzne.

    Teraz už len mocnosti. Vrcholy ležia na dotyčniciach príslušných kružníc, takže
    $$
    \gather
    p(B,\omega_1) = (t-b)^2 = p(B,\omega_2),\\
    p(E,\omega_1) = (e+t)^2 = p(E,\omega_2).
    \endgather
    $$
    Rôzne body $B$, $E$ teda ležia na chordále nesústredných kružníc $\omega_1$, $\omega_2$, tou je preto priamka $BE$. Analogicky hodnoty $(c+t)^2$ v~$C$ a~$(t-f)^2$ v~$F$ robia $CF$ chordálou kružníc $\omega_2$, $\omega_3$ a~hodnoty $(a+t)^2$ v~$A$ a~$(t-d)^2$ v~$D$ robia $AD$ chordálou kružníc $\omega_1$, $\omega_3$.

    Tri po dvoch nesústredné kružnice majú buď potenčný stred, ktorým prechádzajú všetky tri chordály, alebo ich stredy ležia na jednej priamke a~potom sú všetky tri chordály na túto priamku kolmé, teda rovnobežné, prípadne splynú. Druhá možnosť odpadá: uhlopriečka $AD$ delí konvexný šesťuholník na $ABCD$ a~$ADEF$, takže $B$ a~$E$ ležia v~opačných polrovinách určených priamkou $AD$ a~úsečka $BE$ ju pretína v~jedinom bode. Uhlopriečky $AD$, $BE$, $CF$ teda prechádzajú jedným bodom, potenčným stredom kružníc $\omega_1$, $\omega_2$, $\omega_3$.
}

\DisplayExerciseSolutions

\DisplayHints

\DisplayProblemSolutions

\bye
